\chapter{电子波包传播与时间结构}\label{chap:prop-measure}

前面已经求出电子离开相互作用区时的量子态。实验中的探测器通常位于更远的位置，电子还要传播一段距离才到达探测面。这段自由传播会改变不同动量分量的相对相位，即使能谱没有变化，电子在时间和空间中的分布也可能改变。要由出射态预言探测图样，我们先求传播后的态，再计算探测结果的概率。

记相互作用结束时刻为 $t_+$，其电子密度矩阵元为 $\rho_{\ell\ell'}^{(+)}$。以下自由传播中的 $t$ 是从 $t_+$ 起算的经过时间。若第 $\ell$ 个通道的自由能量为 $\mathcal E_\ell$，态矢振幅获得因子 $\exp(-\ii\mathcal E_\ell t/\hbar)$；密度矩阵由一个振幅和另一个振幅的共轭组成，因而
\begin{equation}
\boxed{
\rho_{\ell\ell'}(t)
=\rho_{\ell\ell'}^{(+)}
\exp\!\left[-\frac{\ii t}{\hbar}
(\mathcal E_\ell-\mathcal E_{\ell'})\right].}
\label{eq:sideband_coherence_free_propagation_master}
\end{equation}
当 $\ell=\ell'$ 时，相位因子等于一，各边带的概率保持不变。当 $\ell\ne\ell'$ 时，相对相位继续演化；在位置表象中把各分量叠加，便会出现随传播距离变化的干涉项。后面讨论的电子脉冲列和时间聚焦都来自这一过程。

传播过程中若还有材料涨落或未探测光子参与，电子就需要用约化密度算符描述。我们将在自由传播之后讨论环境怎样改变相干和涨落，最后引入测量算符，说明能谱、相位读出和条件测量分别能够获得哪些信息。这样，同一个输出信号中的传播变化、环境作用和探测器响应就都有明确的计算位置。

\section{自由传播与相空间}

自由粒子的每个动量分量只积累相位，因此在动量表象中传播最容易计算。位置密度则由所有分量的干涉决定，需要在传播后作 Fourier 叠加。本节先写出这一解，再对窄动量分布展开色散，从而得到波包展宽、纵向聚束和时间透镜的计算公式。

\subsection{自由传播与色散曲率矩阵}

设电子在自由传播开始时的动量振幅为 $f_s(\mathbf p)$，其中 $s$ 是自旋标签，归一化为 $\sum_s\int\dd^3p\,|f_s(\mathbf p)|^2=1$。由能量本征态的时间演化，正能电子波包在时刻 $t$ 为
\begin{equation*}
|\Psi(t)\rangle
=\sum_s\int\dd^3 p\,f_s(\mathbf p)
\ee^{-\ii E(\mathbf p)t/\hbar}|\mathbf p,s\rangle,
\qquad
E(\mathbf p)=\sqrt{m_{\mathrm e}^2c^4+c^2p^2}.
\end{equation*}
当实际谱支持足够窄时，局域传播方程由这一精确相位的受控展开得到。

设中心动量 $\mathbf p_0=p_0\mathbf e_z$，记 $\delta\mathbf p=\mathbf p-\mathbf p_0$。对 $E(\mathbf p)$ 作 Taylor 展开。一次导数给出波包中心的群速度；二次导数组成的矩阵称为 Hessian 矩阵，它描述不同动量分量的速度怎样偏离中心速度。保留到二阶，得
\begin{equation*}
E_2(\mathbf p)
=E_0+\mathbf v_0\cdot\delta\mathbf p
+\frac{\delta p_x^2+\delta p_y^2}{2m_\perp}
+\frac{\delta p_z^2}{2m_\parallel},
\end{equation*}
其中
\begin{equation*}
\mathbf v_0=\frac{c^2\mathbf p_0}{E_0},
\qquad
m_\perp=\frac{E_0}{c^2}=\gamma_0m_{\mathrm e},
\qquad
m_\parallel=\gamma_0^3m_{\mathrm e}.
\end{equation*}
$m_\perp$ 与 $m_\parallel$ 分别是同一相对论色散沿横向和纵向方向的局部曲率尺度。传播误差由被舍弃的谱相位而不是单独的带宽控制。对实际动量支持 $\mathcal W_p$ 定义
\begin{equation}
\boxed{
\epsilon_{\rm disp}^{(\ge3)}(t)
\equiv
\sup_{\mathbf p\in\mathcal W_p}
\left|\frac{t}{\hbar}[E(\mathbf p)-E_2(\mathbf p)]\right|.}
\label{eq:dispersion_truncation_control}
\end{equation}
只有在 $\epsilon_{\rm disp}^{(\ge3)}\ll1$ 时，进入共动坐标 $\zeta=z-v_0t$ 后才得到二阶局域传播方程
\begin{equation}
\boxed{
\ii\hbar\partial_t\psi_2
=-\frac{\hbar^2}{2m_\perp}\nabla_\perp^2\psi_2
-\frac{\hbar^2}{2m_\parallel}\partial_\zeta^2\psi_2.}
\label{eq:quadratic_dispersion_local_propagation}
\end{equation}
令 $P_{\mathcal W_p}$ 投影到实际谱支持，精确与二阶传播子的差满足
\begin{equation}
\boxed{
\left\|[U_{\rm ex}(t)-U_2(t)]P_{\mathcal W_p}\right\|
\le\epsilon_{\rm disp}^{(\ge3)}(t).}
\label{eq:quadratic_dispersion_propagator_error_bound}
\end{equation}
该近似的误差由高阶色散相位随传播时间的累积控制，因此必须用无量纲相位而非孤立的哈密顿量系数判断失效。

\begin{proposition}[实谱乘法传播子的二次截断界]
设实色散 $E\in C^3$，归一化波包支撑于有界凸动量窗 $\mathcal W_p$，$E_2$ 为中心处二次 Taylor 多项式。则 $\|U_{\rm ex}-U_2\|_{\mathcal W_p}\le\min\{2,|t|\sup_{\mathcal W_p}|E-E_2|/\hbar\}$；若 $\|D^3E\|\le M_3$，后一个量不超过 $|t|M_3\sup|\delta\mathbf p|^3/(6\hbar)$。对相对论色散，二次包络满足\eqnrefs{eq:quadratic_dispersion_local_propagation}。
\end{proposition}
\begin{proof}
其适用条件是初态动量主要支撑在给定带宽内、$E(p)$ 在该区间足够光滑，并且显式传播相位误差小于一。 Taylor 展开写成

\begin{align*}
E(\mathbf p_0+\delta\mathbf p)
={}&E_0+DE(\mathbf p_0)[\delta\mathbf p]
+\frac12D^2E(\mathbf p_0)[\delta\mathbf p^{\otimes2}]\\
&+\frac1{3!}D^3E(\mathbf p_0)[\delta\mathbf p^{\otimes3}]+R_4.
\end{align*}
对 $E=\sqrt{m_{\mathrm e}^2c^4+c^2p^2}$，
\begin{equation*}
\partial_{p_i}E=\frac{c^2p_i}{E},
\qquad
\partial_{p_i}\partial_{p_j}E
=\frac{c^2}{E}\delta_{ij}-\frac{c^4p_ip_j}{E^3}.
\end{equation*}
在 $\mathbf p_0=p_0\mathbf e_z$ 上，横向 Hessian 元素为 $1/m_\perp=c^2/E_0$，纵向元素为
\begin{align*}
\frac1{m_\parallel}
&=\frac{c^2}{E_0}-\frac{c^4p_0^2}{E_0^3}\\
&=\frac{c^2}{E_0}\left(1-\frac{v_0^2}{c^2}\right)
=\frac{1}{\gamma_0^3m_{\mathrm e}}.
\end{align*}
把载波和群延迟从精确 Fourier 波包中提出。令 $\zeta=z-v_0t$，并写
\begin{equation*}
\Psi(\mathbf r,t)
=\ee^{\ii(p_0z-E_0t)/\hbar}
\psi_2(\mathbf r_\perp,\zeta,t).
\end{equation*}
在二阶色散下，包络为
\begin{equation*}
\psi_2
=\int\frac{\dd^3(\delta p)}{(2\pi\hbar)^3}
A(\delta\mathbf p)
\exp\!\left[
\frac{\ii}{\hbar}(\delta\mathbf p_\perp\cdot\mathbf r_\perp+\delta p_z\zeta)
-\frac{\ii t}{\hbar}
\left(\frac{\delta p_\perp^2}{2m_\perp}
+\frac{\delta p_z^2}{2m_\parallel}\right)\right].
\end{equation*}
对 $t$ 求导，群速度项已经被 $\zeta=z-v_0t$ 吸收，因此
\begin{align*}
\ii\hbar\partial_t\psi_2
&=\int\frac{\dd^3(\delta p)}{(2\pi\hbar)^3}
A(\delta\mathbf p)
\left(\frac{\delta p_\perp^2}{2m_\perp}
+\frac{\delta p_z^2}{2m_\parallel}\right)\ee^{\cdots}\\
&=-\frac{\hbar^2}{2m_\perp}\nabla_\perp^2\psi_2
-\frac{\hbar^2}{2m_\parallel}\partial_\zeta^2\psi_2,
\end{align*}
其中第二行使用 $\delta p_j \ee^{\ii\delta p_jx_j/\hbar}=-\ii\hbar\partial_{x_j}\ee^{\ii\delta p_jx_j/\hbar}$。这就是\eqnrefs{eq:quadratic_dispersion_local_propagation}。精确传播子与二阶传播子在动量表象都是乘法算符。定义余项相位
\begin{equation*}
\delta\varphi(\mathbf p,t)
=\frac{t}{\hbar}[E(\mathbf p)-E_2(\mathbf p)].
\end{equation*}
对任意归一化且支撑在 $\mathcal W_p$ 的波包，
\begin{align*}
\|[U_{\rm ex}(t)-U_2(t)]\psi\|^2
&=\int_{\mathcal W_p}\dd^3 p\,|A(\mathbf p)|^2
|\ee^{-\ii\delta\varphi}-1|^2\\
&\le\sup_{\mathbf p\in\mathcal W_p}|\ee^{-\ii\delta\varphi}-1|^2\\
&\le\sup_{\mathbf p\in\mathcal W_p}|\delta\varphi|^2.
\end{align*}
取平方根后得到\eqnrefs{eq:quadratic_dispersion_propagator_error_bound}。若已知三阶导数上界 $M_3=\sup_{\mathcal W_p}\|D^3E\|$，Taylor 余项还给出
\begin{equation*}
\epsilon_{\rm disp}^{(\ge3)}(t)
\le\frac{|t|M_3}{6\hbar}
\sup_{\mathcal W_p}|\delta\mathbf p|^3+O(|\delta\mathbf p|^4),
\end{equation*}
从而把“带宽小”转化为传播时间相关的相位误差界。
\end{proof}


固定总能量、以 $z$ 作为传播坐标时，同一个质量壳给出
\begin{equation*}
p_z(\mathbf p_\perp)=\sqrt{p_0^2-p_\perp^2}.
\end{equation*}
近轴条件首先要求 $p_\perp^2/p_0^2\ll1$，但长距离传播还要求四阶及以上纵向相位不可分辨。直接定义
\begin{equation}
\boxed{
\epsilon_{\rm par}^{(\ge4)}(z)
\equiv
\sup_{\mathbf p_\perp\in\mathcal W_\perp}
\left|\frac{z}{\hbar}
\left[\sqrt{p_0^2-p_\perp^2}-p_0+\frac{p_\perp^2}{2p_0}\right]\right|.}
\label{eq:paraxial_phase_control_parameter}
\end{equation}
在 $\epsilon_{\rm par}^{(\ge4)}\ll1$ 时，抽去载波 $\ee^{\ii p_0z/\hbar}$ 得到
\begin{equation}
\boxed{
\ii\hbar\partial_z\psi_{\rm par}
=-\frac{\hbar^2}{2p_0}\nabla_\perp^2\psi_{\rm par}.}
\label{eq:relativistic_paraxial_from_mass_shell}
\end{equation}
并有传播子误差界
\begin{equation}
\boxed{
\left\|[U_z^{\rm ex}-U_z^{\rm par}]P_{\mathcal W_\perp}\right\|
\le\epsilon_{\rm par}^{(\ge4)}(z).}
\label{eq:paraxial_propagator_error_bound}
\end{equation}
沿中心轨迹取 $z\simeq v_0t$ 时，$p_0/v_0=\gamma_0m_{\mathrm e}=m_\perp$，因此时间域横向 Hessian 矩阵与定态近轴方程是同一精确色散在不同边值问题下的二阶极限。

\begin{proposition}[正根角谱与抛物传播误差]
固定 $p_0>0$，设角谱支撑于 $|\mathbf p_\perp|\le a<p_0$，取 $p_z=\sqrt{p_0^2-p_\perp^2}$。精确正向传播和二次抛物传播之差满足\eqnrefs{eq:paraxial_propagator_error_bound}，且主导相位余项为 $-zp_\perp^4/(8\hbar p_0^3)$，高阶余项在此窗上一致有界。
\end{proposition}
\begin{proof}
设电子束以 $+z$ 为中心传播方向，并且横向动量与纵向失谐都小于中心动量。固定总能量时有 $p_z^2+p_\perp^2=p_0^2$；取沿 $+z$ 传播的正根，并把高阶角谱项保留为显式余项，

\begin{align*}
p_z
&=p_0\sqrt{1-\frac{p_\perp^2}{p_0^2}}\\
&=p_0\left[1-\frac12\frac{p_\perp^2}{p_0^2}
-\frac18\frac{p_\perp^4}{p_0^4}
+O\!\left(\frac{p_\perp^6}{p_0^6}\right)\right]\\
&=p_0-\frac{p_\perp^2}{2p_0}
-\frac{p_\perp^4}{8p_0^3}+O(p_\perp^6/p_0^5).
\end{align*}
把共同纵向载波抽出，写
\begin{equation*}
\Psi(\mathbf r_\perp,z)
=\ee^{\ii p_0z/\hbar}\psi_{\rm env}(\mathbf r_\perp,z),
\end{equation*}
其中包络的精确角谱表示为
\begin{equation*}
\psi_{\rm env}(\mathbf r_\perp,z)
=\int\frac{\dd^2 p_\perp}{(2\pi\hbar)^2}
\,A(\mathbf p_\perp)
\exp\!\left[\frac{\ii}{\hbar}
\mathbf p_\perp\!\cdot\!\mathbf r_\perp
+\frac{\ii}{\hbar}(p_z-p_0)z\right].
\end{equation*}
对 $z$ 求导时每个角谱分量给出
\begin{align*}
\ii\hbar\,\partial_z\psi_{\rm env}
&=\int\frac{\dd^2 p_\perp}{(2\pi\hbar)^2}
A(\mathbf p_\perp)
\bigl[p_0-p_z(\mathbf p_\perp)\bigr]
\ee^{\ii\mathbf p_\perp\cdot\mathbf r_\perp/\hbar}
\ee^{\ii(p_z-p_0)z/\hbar}\\
&=-\frac{\hbar^2}{2p_0}\nabla_\perp^2\psi_{\rm env}
+\mathcal R_{\rm par}(\mathbf r_\perp,z),
\end{align*}
其中残差被精确定义为
\begin{equation*}
\mathcal R_{\rm par}
\equiv
\int\frac{\dd^2 p_\perp}{(2\pi\hbar)^2}
A(\mathbf p_\perp)
\left[p_0-p_z(\mathbf p_\perp)-\frac{p_\perp^2}{2p_0}\right]
\ee^{\ii\mathbf p_\perp\cdot\mathbf r_\perp/\hbar}
\ee^{\ii(p_z-p_0)z/\hbar}.
\end{equation*}
当 $\mathcal R_{\rm par}$ 相对于保留项低于目标精度时，得到抛物传播方程\eqnrefs{eq:relativistic_paraxial_from_mass_shell}。其适用范围可以直接由精确传播子与抛物传播子的范数差量化。定义
\begin{equation*}
\delta\varphi(\mathbf p_\perp;z)
=\frac{z}{\hbar}
\left[p_z-p_0+\frac{p_\perp^2}{2p_0}\right].
\end{equation*}
两个传播子在角谱表象都为乘法算符，因此对任意支撑在 $\mathcal W_\perp$ 的归一化态，
\begin{align*}
\|[U_z^{\rm ex}-U_z^{\rm par}]\psi\|^2
&=\int_{\mathcal W_\perp}\dd^2 p_\perp\,
|A(\mathbf p_\perp)|^2
|\ee^{\ii\delta\varphi}-1|^2\\
&\le \sup_{\mathcal W_\perp}|\ee^{\ii\delta\varphi}-1|^2
\le \sup_{\mathcal W_\perp}|\delta\varphi|^2,
\end{align*}
其中用了 $|\ee^{\ii x}-1|=2|\sin(x/2)|\le|x|$。取平方根即得到\eqnrefs{eq:paraxial_propagator_error_bound}。由上面的展开还可读出主导误差
\begin{equation*}
\delta\varphi
=-\frac{z\,p_\perp^4}{8\hbar p_0^3}
+O\!\left(\frac{z p_\perp^6}{\hbar p_0^5}\right),
\end{equation*}
所以几何小角条件 $p_\perp/p_0\ll1$ 与长距离相位条件 $z p_\perp^4/(8\hbar p_0^3)\ll1$ 必须同时满足。
\end{proof}


\subsection{能量格相干到时间密度的定量映射}\label{subsec:sideband-to-temporal-density-dp46}

\begin{exbox}[固定探测面的能量相位与传播误差]
时间演化使用 $E(p)t$，固定平面读出则更方便使用 $p(E)L$。
沿向前正能支路，前章的能量归一化满足 $\int|b(E)|^2\dd E=1$。
定义该支路的理想到达振幅
\begin{equation}
 f_L(t)=\int_{m_{\rm e}c^2}^\infty\frac{\dd E}{\sqrt{2\pi\hbar}}\,
 b(E)\exp[\ii p(E)L/\hbar-\ii Et/\hbar],\qquad
 p(E)=\frac{\sqrt{E^2-m_{\rm e}^2c^4}}c.
 \label{eq:feqo-fixed-plane-energy-amplitude}
\end{equation}
Parseval 给 $\int_{\mathbb R}|f_L(t)|^2\dd t=1$，漂移只给能量振幅乘单位模相位。
这个向前支路的理想时间测量是一个 POVM；它不是对任意反射波都成立的“位置密度等于到达概率”。
真实探测器还需使用后文的效应和响应核。

为求群延迟，先直接微分而不在 $p$、波数和能量之间混用单位：
\begin{equation}
 p'(E)=\frac{E}{c\sqrt{E^2-m_{\rm e}^2c^4}}=\frac1v,\qquad
 p''(E)=-\frac{m_{\rm e}^2c^3}{(E^2-m_{\rm e}^2c^4)^{3/2}}
 =-\frac1{\gamma^3m_{\rm e}v^3}.
 \label{eq:feqo-fixed-plane-momentum-derivatives}
\end{equation}
这里 $p=\gamma m_{\rm e}v$，故第二个等式立即由
$(E^2-m_{\rm e}^2c^4)^{3/2}=c^3p^3$ 得到。
记 $\varepsilon=E-E_0$、$\tau=t-L/v_0$，除去共同载波后的振幅为
\begin{equation}
 \widetilde f_L(\tau)=\int\frac{\dd\varepsilon}{\sqrt{2\pi\hbar}}\,
 b(E_0+\varepsilon)\exp\!\left[-\frac{\ii\varepsilon\tau}{\hbar}
 -\frac{\ii L\varepsilon^2}{2\hbar\gamma_0^3m_{\rm e}v_0^3}
 +\frac{\ii L R_3(\varepsilon)}\hbar\right].
 \label{eq:feqo-fixed-plane-fresnel-kernel}
\end{equation}
其中 $|R_3|\le M_{p,3}|\varepsilon|^3/6$。
在实际支持 $|\varepsilon|\le\Delta E$ 上，丢弃三阶相位的 $L^2(\dd t)$ 振幅误差至多为
$LM_{p,3}(\Delta E)^3/(6\hbar)$；无有限支集的包还需计入尾部。
这里直接采用能量变量，不需要再把传播时间粗略替换为 $L/v_0$。
若从二阶动量 Fourier 积分的 $t=L/v_0+\tau$ 出发，另行丢弃 $\tau$ 乘二次色散的项，
还必须检查 $|\tau|\Delta p^2/(2m_\parallel\hbar)\ll1$；仅有 $v_0|\tau|\ll L$ 并不能保证该累计相位很小。

相位的驻点给局部群延迟
\begin{equation}
 t(E)=Lp'(E)=\frac L{v(E)},\qquad
 \delta t=-\frac{L\delta E}{\gamma_0^3m_{\rm e}v_0^3}+O(\delta E^2).
 \label{eq:feqo-arrival-time-shift}
\end{equation}
电子获得能量时到达更早。相邻能量梳分量的中心延迟差领先阶为
$-L\hbar\omega/(\gamma_0^3m_{\rm e}v_0^3)$。
相干源包之间是否还能干涉，需将这个分离量与源包宽度比较；
逐发中心能量随机平均的相位衰减则是另一统计操作，不能把两者统称为同一内禀退相干长度。
\end{exbox}

\begin{figure}[htbp]
 \centering
 \includegraphics[width=0.55\linewidth]{pinem_energy_ladder.pdf}
 \caption{正负整数边带的能量标记 $E_0+\ell\hbar\omega$。这张图表示能量阶数，
 不表示边带占据、真实动量等间隔或任意能区都可延伸到负无限的物理谱。}
 \label{fig:pinem-energy-ladder}
\end{figure}


前面的自由传播一般表达式还可以直接化成实验可读取的纵向密度。对\chapref{chap:external-pinem}定义的单频交换格，在中心同步条件 $v_0\kappa_z=\omega$ 下，相邻电子通道在共动坐标 $\zeta=z-v_0t$ 中相差纵向相位 $\kappa_z\zeta$。若只保留二阶纵向色散，则传播后的约化密度核给出
\begin{equation}
\boxed{
 \mathfrak n_e(\zeta,t)
 =\sum_{\ell,\ell'}
 \rho_{\ell\ell'}^{(+)}
 \exp\!\left[\ii(\ell-\ell')\kappa_z\zeta
 -\ii(\ell^2-\ell'^2)\omega_{\rm r}t\right].}
\label{eq:sideband_to_longitudinal_density_dp46}
\end{equation}
这里 $\mathfrak n_e$ 是把慢变纵向包络的共同归一化因子抽出后得到的无量纲周期密度；若需要绝对线密度，只需乘入入射包络线密度。\eqnrefs{eq:sideband_to_longitudinal_density_dp46} 显示：对角能谱 $\rho_{\ell\ell}$ 只给常数背景，所有时间结构都来自能量格非对角相干。

把谐波阶数定义为 $\Delta\ell=\ell-\ell'$，可将密度写成
\begin{equation}
 \mathfrak n_e(\zeta,t)
 =\sum_{\Delta\ell\in\mathbb Z}B_{\Delta\ell}(t)\ee^{\ii \Delta\ell\kappa_z\zeta},
 \qquad B_{-\Delta\ell}=B_{\Delta\ell}^*,
\label{eq:electron_density_harmonic_expansion_dp46}
\end{equation}
其中
\begin{equation}
\boxed{
 B_{\Delta\ell}(t)
 =\sum_{\ell}
 \rho_{\ell,\ell-\Delta\ell}^{(+)}
 \exp\!\left[-\ii(2\ell \Delta\ell-\Delta\ell^2)\omega_{\rm r}t\right].}
\label{eq:electron_bunching_harmonic_dp46}
\end{equation}
$B_{\Delta\ell}$ 是第 $\Delta\ell$ 阶电子聚束因子，也就是归一化密度在该空间频率上的复 Fourier 系数。$|B_{\Delta\ell}|$ 给出这一谐波的幅度，$\arg B_{\Delta\ell}$ 给出它相对于光学相位基准的位置；若密度只含常数项与一对 $\pm\Delta\ell$ 谐波，则对应余弦项的峰值幅度为 $2|B_{\Delta\ell}|$。尤其
\begin{equation*}
 B_0=\sum_\ell\rho_{\ell\ell}^{(+)}=1,
\end{equation*}
因此归一化由密度矩阵迹自动保证。

二阶反冲相位是 $\ell^2\omega_{\rm r}t$，故所有整数格点相位在
\begin{equation}
\boxed{
 T_{\rm T}=\frac{2\pi}{\omega_{\rm r}},
 \qquad
 L_{\rm T}=v_0T_{\rm T}
 =\frac{2\pi v_0}{\omega_{\rm r}}}
\label{eq:electron_talbot_time_length_dp46}
\end{equation}
后同时复现。$T_{\rm T}$ 与 $L_{\rm T}$ 分别是电子能量格的 Talbot 时间和沿中心轨迹的 Talbot 长度；二者均由同一个纵向色散质量
$m_\parallel=\gamma_0^3m_{\mathrm e}$ 和动量步长 $\hbar \kappa_z$ 决定：
\begin{equation*}
 L_{\rm T}
 =\frac{4\pi m_\parallel v_0}{\hbar \kappa_z^2}.
\end{equation*}
若使用同步关系 $\kappa_z=\omega/v_0$，还可写成
\begin{equation*}
 L_{\rm T}
 =\frac{4\pi m_\parallel v_0^3}{\hbar\omega^2}.
\end{equation*}
Talbot 复现还可以直接从整数格的数论相位看出。半个 Talbot 时间 $t=T_{\rm T}/2$ 时，
\begin{equation*}
 \ee^{-\ii\ell^2\omega_{\rm r}T_{\rm T}/2}
 =\ee^{-\ii\pi\ell^2}
 =\ee^{-\ii\pi\ell},
\end{equation*}
因为整数满足 $\ell^2-\ell=\ell(\ell-1)\in2\mathbb Z$。因此每个格点多出的因子 $\ee^{-\ii\pi\ell}$ 可吸收到空间相位中：
\begin{equation*}
 \ee^{\ii\ell\kappa_z\zeta}\ee^{-\ii\pi\ell}
 =\ee^{\ii\ell\kappa_z(\zeta-\pi/\kappa_z)}.
\end{equation*}
所以半 Talbot 时间并非原位复现，而是把整个周期密度平移半个空间周期 $\pi/\kappa_z$；到 $T_{\rm T}$ 时该平移再消失，得到原位复现。这个结果来自离散二次谱相位本身，与 PINEM 初始边带振幅的具体形状无关。

\begin{proposition}[整数能量梳的二次相位复现]
对整数边带 $\ell\in\mathbb Z$，二次传播相位 $e^{-\ii s\ell^2}$ 在 $s=2\pi$ 严格复现，在 $s=\pi$ 等价于 $e^{-\ii\pi\ell}$ 的半周期平移。二次色散系数换算为实验时间与长度后给\eqnrefs{eq:electron_talbot_time_length_dp46}；高阶色散存在时复现仅为受控近似。
\end{proposition}
\begin{proof}
设占据的能量边带由整数 $\ell\in\mathbb Z$ 标记，并先只保留二次色散。第 $\ell$ 个边带在传播长度 $L_z$ 后获得
\[
U_\ell^{(2)}(L_z)
=\exp[-\ii\ell^2\vartheta_2(L_z)],
\qquad
\vartheta_2(L_z)=\frac{\hbar\omega^2L_z}{2m_\parallel v_0^3}.
\]
若传播前纯态写成
\[
|\psi(0)\rangle=\sum_\ell c_\ell|\ell\rangle,
\]
传播后为
\[
|\psi(L_z)\rangle
=\sum_\ell c_\ell
\ee^{-\ii\ell^2\vartheta_2(L_z)}|\ell\rangle.
\]
密度矩阵元因此满足
\begin{align*}
\rho_{\ell m}(L_z)
&=c_\ell c_m^*
\ee^{-\ii\ell^2\vartheta_2}
\ee^{+\ii m^2\vartheta_2}\\
&=\rho_{\ell m}(0)
\ee^{-\ii(\ell^2-m^2)\vartheta_2}.
\end{align*}
对混态同样成立，因为自由传播在线性算符层面逐个矩阵元乘上同一相对相位。

完全原位复现要求所有已占据矩阵元同时恢复。若希望得到与具体占据集合无关的整数格复现条件，就要求
\[
\ee^{-\ii(\ell^2-m^2)\vartheta_2(L_T)}=1
\qquad
\forall\,\ell,m\in\mathbb Z.
\]
取 $m=0$，必要条件为
\[
\ell^2\vartheta_2(L_T)=2\pi n_\ell,
\qquad n_\ell\in\mathbb Z.
\]
其中 $\ell=1$ 已经要求
\[
\vartheta_2(L_T)=2\pi n_1.
\]
最小正复现取 $n_1=1$，即
\[
\vartheta_2(L_T)=2\pi.
\]
反过来，把这个值代回任意整数 $\ell,m$：
\begin{align*}
(\ell^2-m^2)\vartheta_2(L_T)
&=2\pi(\ell^2-m^2),
\end{align*}
右端总是 $2\pi$ 的整数倍，所以该条件不仅必要，而且对完整整数能量格充分。于是
\begin{align*}
2\pi
&=\frac{\hbar\omega^2L_T}{2m_\parallel v_0^3},\\
L_T
&=\frac{4\pi m_\parallel v_0^3}{\hbar\omega^2},\\
T_T=\frac{L_T}{v_0}
&=\frac{4\pi m_\parallel v_0^2}{\hbar\omega^2},
\end{align*}
得到正文\eqnrefs{eq:electron_talbot_time_length_dp46}。

半个 Talbot 长度对应
\[
\vartheta_2(L_T/2)=\pi.
\]
此时单边带相位为
\begin{align*}
U_\ell^{(2)}(L_T/2)
&=\ee^{-\ii\pi\ell^2}\\
&=(-1)^{\ell^2}\\
&=(-1)^\ell,
\end{align*}
最后一步使用整数 $\ell$ 与 $\ell^2$ 具有相同奇偶性。另一方面，共动坐标平移
\[
\Delta\zeta=\frac{\pi}{\kappa_z}
\]
会让第 $\ell$ 个空间 Fourier 分量获得
\[
\ee^{\ii\ell\kappa_z\Delta\zeta}
=\ee^{\ii\pi\ell}=(-1)^\ell.
\]
因此 $L_T/2$ 的传播与半周期平移具有相同的格点相位作用，而 $L_T$ 才对应原位完全复现。
\end{proof}


实际电子只占据有限边带，并且精确相对论色散还含三阶及更高阶谱相位。若有效占据范围为 $|\ell|\le\ell_{\max}$，把一个 Talbot 长度内所有三阶以上附加相位的最大值定义为
\begin{equation}
 \boxed{
 \epsilon_T^{(\ge3)}
 \equiv
 \max_{|\ell|\le\ell_{\max}}
 \left|\varphi_\ell^{(\ge3)}(L_T)\right|.}
 \label{eq:electron-talbot-higher-order-control-dp76}
\end{equation}
只有 $\epsilon_T^{(\ge3)}\ll1$ 时，二次谱相位给出的整数复现才保持高对比度；否则必须回到完整相对论传播相位逐个边带计算。

由 $m_\parallel=\gamma_0^3m_{\mathrm e}$ 可直接读出：提高中心 Lorentz 因子会增大纵向色散质量，因而在固定 $\kappa_z$ 下减小 $\omega_{\rm r}=\hbar\kappa_z^2/(2m_\parallel)$；形成同一二次谱相位所需的漂移时间与距离随之增加。

对理想经典 PINEM 输出
\begin{equation*}
 c_\ell^{(+)}
 =J_\ell(2|\beta|)
 \ee^{\ii\ell\varphi_\beta},
 \qquad
 \rho_{\ell\ell'}^{(+)}=c_\ell^{(+)}c_{\ell'}^{(+)*},
\end{equation*}
在理想双边 PINEM 输出并采用前述二阶纵向色散时，\eqnrefs{eq:electron_bunching_harmonic_dp46} 变为完全可计算的闭式级数
\begin{equation}
\boxed{
 B_{\Delta\ell}^{\rm PINEM}(t)
 =\ee^{\ii \Delta\ell\varphi_\beta}
 \sum_{\ell=-\infty}^{\infty}
 J_\ell(2|\beta|)J_{\ell-\Delta\ell}(2|\beta|)
 \ee^{-\ii(2\ell \Delta\ell-\Delta\ell^2)\omega_{\rm r}t}.}
\label{eq:pinem_bunching_exact_series_dp46}
\end{equation}
它给出“耦合强度 $|\beta|$、漂移距离 $L\simeq v_0t$、反冲频率 $\omega_{\rm r}$”到任意密度谐波的直接前向映射，不再需要用“压缩点”“最佳聚束”作定性判断。若实验只读取第一谐波，可先定义第一谐波的相对调制幅度
\begin{equation}
\boxed{\mathcal M_1(t)=2|B_1(t)|,}
\label{eq:first_harmonic_visibility_dp46}
\end{equation}
它一般不等于限制在 $[0,1]$ 的 Michelson 可见度；只有高阶谐波可忽略时二者才一致。完整密度对比度应直接由\eqnrefs{eq:sideband_to_longitudinal_density_dp46} 的最大、最小值计算。

\begin{proposition}[周期边带态的密度 Fourier 系数]
设归一化周期包络具有 Fourier 系数 $c_\ell$，$\sum_\ell|c_\ell|^2=1$。其密度的第 $d$ 个 Fourier 系数为 $\sum_\ell c_\ell c_{\ell-d}^*$，该和由 Cauchy--Schwarz 绝对收敛。代入纯正弦相位产生的 Bessel 系数及二次自由传播相位，得到\eqnrefs{eq:pinem_bunching_exact_series_dp46}。
\end{proposition}
\begin{proof}
先把共同的慢变纵向包络抽出，只处理一个同步周期内的离散能量格。取
\[
L_\zeta\equiv\frac{2\pi}{\kappa_z},
\qquad
\langle\zeta|\ell\rangle
=\frac{1}{\sqrt{L_\zeta}}
\ee^{+\ii\ell\kappa_z\zeta},
\qquad 0\le\zeta<L_\zeta,
\]
则 $\{|\ell\rangle\}_{\ell\in\mathbb Z}$ 在该周期盒中正交归一：
\begin{align*}
\langle\ell'|\ell\rangle
&=\frac1{L_\zeta}
\int_0^{L_\zeta}\dd\zeta\,
\ee^{+\ii(\ell-\ell')\kappa_z\zeta}\\
&=\delta_{\ell\ell'}.
\end{align*}
盒归一化的位置密度为
$\langle\zeta|\rho(t)|\zeta\rangle$；正文使用的无量纲周期密度定义为
\[
\mathfrak n_e(\zeta,t)
\equiv L_\zeta\langle\zeta|\rho(t)|\zeta\rangle,
\]
所以共同的 $L_\zeta^{-1}$ 因子不会进入最终谐波系数。

在二阶纵向色散下，共动表象已经移除了 $E_0+\ell\hbar\omega$ 的常数与线性相位，只留下反冲项 $\hbar\omega_{\rm r}\ell^2$。因此正文\eqnrefs{eq:sideband_coherence_free_propagation_master} 给
\[
\rho_{\ell\ell'}(t)
=\rho_{\ell\ell'}^{(+)}
\ee^{-\ii(\ell^2-\ell'^2)\omega_{\rm r}t}.
\]
在位置密度中插入两次能量格完备关系：
\begin{align*}
\mathfrak n_e(\zeta,t)
&=L_\zeta
\sum_{\ell,\ell'}
\langle\zeta|\ell\rangle
\rho_{\ell\ell'}(t)
\langle\ell'|\zeta\rangle\\
&=\sum_{\ell,\ell'}
\rho_{\ell\ell'}^{(+)}
\ee^{+\ii(\ell-\ell')\kappa_z\zeta}
\ee^{-\ii(\ell^2-\ell'^2)\omega_{\rm r}t},
\end{align*}
这就是正文\eqnrefs{eq:sideband_to_longitudinal_density_dp46}。

为了把双重和式按可测空间谐波重新分组，定义
\[
r\equiv\Delta\ell=\ell-\ell',
\qquad
\ell'=\ell-r.
\]
二次相位差逐项化为
\begin{align*}
\ell^2-\ell'^2
&=\ell^2-(\ell-r)^2\\
&=\ell^2-(\ell^2-2\ell r+r^2)\\
&=2\ell r-r^2.
\end{align*}
因此
\begin{align*}
\mathfrak n_e(\zeta,t)
&=\sum_{r\in\mathbb Z}
\ee^{+\ii r\kappa_z\zeta}
\sum_{\ell\in\mathbb Z}
\rho_{\ell,\ell-r}^{(+)}
\ee^{-\ii(2\ell r-r^2)\omega_{\rm r}t}\\
&\equiv
\sum_{r\in\mathbb Z}B_r(t)
\ee^{+\ii r\kappa_z\zeta},
\end{align*}
从而
\[
B_r(t)
=\sum_\ell
\rho_{\ell,\ell-r}^{(+)}
\ee^{-\ii(2\ell r-r^2)\omega_{\rm r}t},
\]
即正文\eqnrefs{eq:electron_bunching_harmonic_dp46}。密度矩阵厄米性
$\rho_{ab}=\rho_{ba}^*$ 再给
\begin{align*}
B_{-r}(t)
&=\sum_\ell
\rho_{\ell,\ell+r}^{(+)}
\ee^{+\ii(2\ell r+r^2)\omega_{\rm r}t}\\
&=\left[
\sum_j
\rho_{j,j-r}^{(+)}
\ee^{-\ii(2jr-r^2)\omega_{\rm r}t}
\right]^*
=B_r^*(t),
\end{align*}
其中第二行取 $j=\ell+r$。所以 $\mathfrak n_e$ 自动为实函数；而
\[
B_0=\sum_\ell\rho_{\ell\ell}^{(+)}
=\operatorname{Tr}\rho^{(+)}=1
\]
保证周期平均密度归一化。

最后代入理想经典 PINEM 纯态振幅
\[
c_\ell^{(+)}
=J_\ell(2|\beta|)\ee^{+\ii\ell\varphi_\beta},
\]
有
\begin{align*}
\rho_{\ell,\ell-r}^{(+)}
&=c_\ell^{(+)}c_{\ell-r}^{(+)*}\\
&=J_\ell(2|\beta|)J_{\ell-r}(2|\beta|)
\ee^{+\ii r\varphi_\beta}.
\end{align*}
代入 $B_r(t)$ 并把与 $\ell$ 无关的相位提出求和号，得到
\[
\boxed{
B_r^{\rm PINEM}(t)
=\ee^{+\ii r\varphi_\beta}
\sum_{\ell=-\infty}^{\infty}
J_\ell(2|\beta|)J_{\ell-r}(2|\beta|)
\ee^{-\ii(2\ell r-r^2)\omega_{\rm r}t},}
\]
即正文\eqnrefs{eq:pinem_bunching_exact_series_dp46}。
\end{proof}


若存在纯退相干而不改变边带占据，可把第 $\Delta\ell$ 条副对角线写成
\begin{equation*}
 \rho_{\ell,\ell-\Delta\ell}^{(+)}
 \longrightarrow
 \mathsf D_{\Delta\ell}^{\rm coh}\rho_{\ell,\ell-\Delta\ell}^{(+)},
 \qquad 0\le|\mathsf D_{\Delta\ell}^{\rm coh}|\le1.
\end{equation*}
于是
\begin{equation}
\boxed{B_{\Delta\ell}(t)\longrightarrow\mathsf D_{\Delta\ell}^{\rm coh} B_{\Delta\ell}(t),}
\label{eq:bunching_dephasing_transfer_dp46}
\end{equation}
而 $B_0$ 不变。这样开放系统中的相干损失直接变成可测密度谐波的乘性衰减，从而把退相干、层析和 Fisher 信息都连接到同一可测密度谐波。

\subsection{从能谱相干算出电子脉冲列}
\label{sec:electron-temporal-bunching-worked}

在能谱上看到一串边带以后，还不能仅凭各峰的高度确定电子的时间分布；不同边带的相位决定它们在何时相长干涉。下面把上一节的密度矩阵传播式用于 PINEM 输出，求出周期密度及其 Fourier 谐波。令 $\theta_\zeta=\kappa_z\zeta$ 表示一个空间周期内的位置，$s_{\rm d}=\omega_{\rm r}t\simeq\omega_{\rm r}L/v_0$ 表示传播中积累的反冲相位；两者均无量纲。相对周期密度为
\begin{equation}
 \mathfrak n_e(\theta_\zeta,s_{\rm d})
 =\sum_{\ell m}\rho_{\ell m}^{(+)}
 \ee^{\ii(\ell-m)\theta_\zeta-\ii(\ell^2-m^2)s_{\rm d}},\qquad
 \frac1{2\pi}\int_0^{2\pi}\mathfrak n_e\dd\theta_\zeta=1.
 \label{eq:electron-bunching-normalized-density}
\end{equation}
$\mathfrak n_e$ 是相对于慢包络背景的无量纲密度，不是单位体积内的电子数。在固定探测面，相对于中心电子的到达延迟为 $\tau_{\rm d}=t-L/v_0$，所以 $\theta_\zeta=-\omega\tau_{\rm d}$；正的到达延迟对应负的共动坐标。把 $s_{\rm d}$ 在一个光学周期内视为常数，还要求实际占据阶数上 $\ell^2\omega_{\rm r}|\tau_{\rm d}|\ll1$。

对任意混态，$\rho_{\ell m}^{(+)}$ 的对角部分给常背景，非对角部分才给随 $\theta_\zeta$ 变化的项。若能谱为 $P_\ell$ 而相位完全随机，则 $\rho_{\ell m}=P_\ell\delta_{\ell m}$，无论传播多远都有 $\mathfrak n_e=1$。相反，纯 PINEM 态的 $c_\ell=J_\ell(2|\beta|)\ee^{\ii\ell\varphi_\beta}$ 保留确定的交叉项。刚离开相互作用区时
\begin{equation*}
 \sum_\ell c_\ell\ee^{\ii\ell\theta_\zeta}
 =\exp[2\ii|\beta|\sin(\theta_\zeta+\varphi_\beta)],
\end{equation*}
模平方仍为一；后续二次色散才把相位调制变为密度调制。产生能谱边带与产生时间脉冲是不同步骤。

将周期密度写为 $\mathfrak n_e=\sum_d B_d\ee^{\ii d\theta_\zeta}$，其 Fourier 系数为
\[
B_d=\frac1{2\pi}\int_0^{2\pi}\dd\theta_\zeta\,
\mathfrak n_e(\theta_\zeta,s_{\rm d})\ee^{-\ii d\theta_\zeta}.
\]
把密度的双重求和代入，角积分只留下 $\ell-m=d$ 的项。这解释了为什么第 $d$ 阶密度谐波需要所有相隔 $d$ 条边带的相干共同贡献。对上述纯 PINEM 输入，Bessel 和式进一步给出
\begin{proposition}[Bessel 梳的完整密度谐波]
\begin{equation}
 B_d(s_{\rm d})=(-\ii)^d\ee^{\ii d\varphi_\beta}
 J_d[4|\beta|\sin(ds_{\rm d})],\qquad d\in\mathbb Z.
 \label{eq:electron-pinem-bunching-closed}\end{equation}
\end{proposition}

\begin{proof}
用 $u(\theta)=\exp(\ii x\sin\theta)=\sum_\ell J_\ell(x)\ee^{\ii\ell\theta}$，
考察 $u(\theta-\alpha/2)u(\theta+\alpha/2)^*$。
直接合并指数给
\[
u(\theta-\alpha/2)u(\theta+\alpha/2)^*
=\exp[-2\ii x\sin(\alpha/2)\cos\theta]
=\sum_d(-\ii)^dJ_d(2x\sin(\alpha/2))\ee^{\ii d\theta}.
\]
另一方面，其第 $d$ 个 Fourier 系数为
$\sum_\ell J_\ell(x)J_{\ell-d}(x)\ee^{-\ii(\ell-d/2)\alpha}$。
取 $x=2|\beta|$、$\alpha=2ds_{\rm d}$，再乘初态因子
$\ee^{\ii d\varphi_\beta}$，即得到闭式，同时固定其复相位。
\end{proof}

\begin{proposition}[Bessel 副对角线和的闭式]
设 $x\in\mathbb R$、$d\in\mathbb Z$。Bessel 副对角线级数绝对收敛，加入二次传播相位后，\eqnrefs{eq:pinem_bunching_exact_series_dp46}可精确求和为\eqnrefs{eq:electron-pinem-bunching-closed}。独立高斯相位抖动另乘其特征函数\eqnrefs{eq:electron-bunching-jitter-transfer}。
\end{proposition}
\begin{proof}
令 $x=2|\beta|$。先把\eqnrefs{eq:pinem_bunching_exact_series_dp46} 中的阶差固定为 $d$，而不截断其中的边带和：
\begin{equation*}
 B_d=\ee^{\ii d\varphi_\beta+\ii d^2s_{\rm d}}
 \sum_{\ell\in\mathbb Z}J_\ell(x)J_{\ell-d}(x)
 \ee^{-2\ii\ell ds_{\rm d}}.
\end{equation*}
对两个实 Bessel 系数分别使用互为共轭的 Fourier 表示，
\begin{equation*}
 J_\ell(x)=\frac1{2\pi}\int_0^{2\pi}\ee^{\ii x\sin u-\ii\ell u}\dd u,
 \qquad
 J_{\ell-d}(x)=\frac1{2\pi}\int_0^{2\pi}
 \ee^{-\ii x\sin v+\ii(\ell-d)v}\dd v.
\end{equation*}
此时所有 $\ell$ 依赖都集中在
$\sum_\ell\exp[-\ii\ell(u-v+2ds_{\rm d})]$。
它是 $2\pi$ 倍的周期 Dirac $\delta$，要求 $v=u+2ds_{\rm d}$ 模 $2\pi$；该等式按周期分布理解，也可先加 $\ee^{-\epsilon\ell^2}$ 收敛因子再令 $\epsilon\to0^+$。因此双角积分只剩一个。使用
\begin{equation*}
 \sin u-\sin(u+2ds_{\rm d})
 =-2\sin(ds_{\rm d})\cos(u+ds_{\rm d}),
\end{equation*}
并把剩余角平移为 $u+ds_{\rm d}$，外部的 $\ee^{\ii d^2s_{\rm d}}$ 恰与平移产生的相位抵消，得到
\begin{equation*}
 B_d=\frac{\ee^{\ii d\varphi_\beta}}{2\pi}
 \int_0^{2\pi}\dd u\,\ee^{-\ii du}
 \exp[-2\ii x\sin(ds_{\rm d})\cos u].
\end{equation*}
使用 $\exp(-\ii y\cos u)=\sum_r(-\ii)^rJ_r(y)\ee^{\ii ru}$，角积分只保留 $r=d$，便得到\eqnrefs{eq:electron-pinem-bunching-closed}。最后对抖动相位的高斯特征函数取平均，给\eqnrefs{eq:electron-bunching-jitter-transfer}。这两个步骤分别使用确定性自由传播和不同实验重复之间的统计平均，不能把后者当作额外的确定性传播相位。
\end{proof}



$d$ 是被读取的密度谐波阶数，不是某一个末态边带编号；$B_d$ 汇总所有相隔 $d$ 阶的相干。该式保持 $B_0=1$、$B_{-d}=B_d^*$，并在 $s_{\rm d}=0$ 时使所有非零阶谐波消失。

\begin{exbox}[第一谐波最强的距离：弱调制与强调制]
给定 $b=|\beta|>0$，求第一谐波首次达到最大值的漂移距离。由闭式，问题化为在 $0\le s_{\rm d}\le\pi/2$ 内最大化
$|B_1|=|J_1(4b\sin s_{\rm d})|$。

\textbf{解。}令 $x_*$ 为 $J_1'(x)=0$ 的第一个正根。递推关系给
$J_1'(x)=[J_0(x)-J_2(x)]/2$，所以该常数也由 $J_0(x_*)=J_2(x_*)$ 定义；$x_*\simeq1.84118$。$J_1$ 的首个正极大值也是其绝对最大值。现在只需检查驱动能否达到这个参数：
\begin{equation}
 s_{\rm d,*}=\begin{cases}
 \pi/2,&4b<x_*,\\
 \arcsin[x_*/(4b)],&4b\ge x_*,
 \end{cases}\qquad
 L_* =\frac{v_0s_{\rm d,*}}{\omega_{\rm r}}.
\label{eq:feqo-optimal-drift}
\end{equation}
弱调制时可用 $J_1(x)\simeq x/2$，故 $|B_1|\simeq2b|\sin s_{\rm d}|$，首个最优位置是 $L_T/4$。强调制时展开 $\arcsin x\simeq x$，得
$L_*\simeq v_0x_*/(4b\omega_{\rm r})$，最优距离随场幅近似反比缩短。

达到 $x_*$ 后，第一谐波上限为 $|B_1|=J_1(x_*)\simeq0.581865$；对应 $\mathcal M_1=2|B_1|>1$。这不违反密度正性，因为 $\mathcal M_1$ 是 Fourier 调制幅度。若只保留 $1+2\Re(B_1\ee^{\ii\theta_\zeta})$，反而会出现负值，说明必须保留更高谐波。第一谐波最大、完整脉冲最窄和 Talbot 复现是三个不同条件。
\end{exbox}

到达时间抖动会改变不同次电子相对于驱动的相位。设逐发误差 $\delta t$ 为零均值高斯变量，标准差为 $\sigma_{t,\rm jit}$，且仅引起相位平移、不改变脉冲包络耦合，则
\begin{equation}
 \overline B_d=B_d\,\langle\ee^{\ii d\omega\delta t}\rangle
 =B_d\exp[-d^2\omega^2\sigma_{t,\rm jit}^2/2].
 \label{eq:electron-bunching-jitter-transfer}
\end{equation}
这就是为何高阶时间结构比总能谱更容易消失：抖动不改变 $P_\ell$，却按 $d^2$ 更快抹去高阶相干。探测器自身的高斯时间响应也给相同形式的谐波衰减，但它只模糊记录，不能据此断言到达探测器前的量子态已经退相干。

\begin{exbox}[从所需谐波反推时间抖动容差]
要求第 $d\ne0$ 阶谐波至少保留理想幅度的比例 $0<\eta_d<1$。求允许的抖动标准差与可用谐波阶数。

\textbf{解。}由高斯衰减因子取对数，直接得到
\begin{equation*}
 \exp[-d^2\omega^2\sigma_{t,\rm jit}^2/2]\ge\eta_d
 \ \Longleftrightarrow\ 
 \sigma_{t,\rm jit}\le\frac{\sqrt{2\ln(1/\eta_d)}}{|d|\omega}.
\end{equation*}
因此分辨越高的谐波，允许的抖动按 $1/|d|$ 缩小。若给定同一保留阈值 $\eta$，可用谐波阶数的上限为
\begin{equation*}
 d_{\max}=\left\lfloor
 \frac{\sqrt{2\ln(1/\eta)}}{\omega\sigma_{t,\rm jit}}
 \right\rfloor,
\end{equation*}
其中向下取整表示取不超过括号内实数的最大整数。

若探测器另有独立高斯时间响应宽度 $\sigma_{t,\rm det}$，两个特征函数相乘，测得谐波为
\begin{equation*}
 B_d^{\rm obs}=B_d\exp\!\left[-\frac{d^2\omega^2}{2}
 (\sigma_{t,\rm jit}^2+\sigma_{t,\rm det}^2)\right].
\end{equation*}
因此仅凭谐波衰减，不能区分束流到达抖动与探测时间模糊。若 $B_d\ne0$，不同 $d$ 的 $-\ln|B_d^{\rm obs}/B_d|/d^2$ 应为同一常数；偏离这个关系说明单一高斯模糊模型不充分。这给出不依赖某组特定时间数值的实验检查。
\end{exbox}

\figref{fig:electron-bunching-worked} 比较相干态、相位混合以及有限抖动后的计算结果。
\begin{figure}[htbp]
 \centering
 \includegraphics[width=0.92\linewidth]{electron_bunching_worked.pdf}
 \caption{能谱到时间密度需要非对角相干。数值图取 $|\beta|=1.20$、$\SI{200}{\kilo\electronvolt}$ 电子及 $\SI{800}{\nano\metre}$ 同步光学步长。左图用相同的 PINEM 占据比较纯态、完全相位混合和有限到达抖动在 $L_*$ 的周期密度；右图给出前三阶聚束因子随漂移的变化。实空间相位与探测面到达延迟满足 $\theta_\zeta=-\omega\tau_{\rm d}$。}
 \label{fig:electron-bunching-worked}
\end{figure}

\begin{figure}[htbp]
 \centering
 \includegraphics[width=0.76\linewidth]{tomography_jitter_visibility.pdf}
 \caption{同一高斯到达抖动对第 $q$ 条副对角相干的抑制，
 $V_q=\exp[-q^2\omega^2\sigma_t^2/2]$，波长 $\SI{800}{\nano\metre}$。
 三种抖动标准差为 $50,100,250$ 阿秒；高次相干比相邻边带相干更早失去可见性。}
 \label{fig:tomography-jitter-visibility}
\end{figure}


\begin{exbox}[三边带截断的归一化与第一谐波]
沿用前面的 $b=|\beta|$，并记 $s=\omega_{\rm r}t=s_{\rm d}$。
三边带近似则保留 $c_0=J_0(2b)$、$c_{\pm1}=\pm J_1(2b)
\ee^{\pm\ii\varphi_\beta}$。其保留概率为
$S_3=J_0(2b)^2+2J_1(2b)^2$，遗漏概率为 $1-S_3$。
若条件于这三个通道，并令 $\theta=\kappa_z\zeta+\varphi_\beta$，周期密度是
\begin{equation}
\mathfrak n_3(\theta,s)=\frac{1}{S_3}
\left[J_0^2+4J_1^2\sin^2\theta
+4J_0J_1\sin s\sin\theta\right].
\label{eq:feqo-three-sideband-density}
\end{equation}
其第一 Fourier 系数的模长为 $2|J_0J_1\sin s|/S_3$，
正弦调制幅度是这个数的两倍。
三边带模型的极大在 $s=\pi/2$，但当遗漏概率不小，它不能替代完整 Bessel 梳的优化。
\end{exbox}

\begin{exbox}[相对论漂移长度与有限能散的数值核对]
取动能 $K_0=\SI{200}{\kilo\electronvolt}$、波长 $\lambda=\SI{800}{\nano\metre}$。
有 $\gamma_0=1.39139$、$v_0=2.08450\times10^8\,\si{\metre\per\second}$，
$\omega=2.35456\times10^{15}\,\si{\per\second}$，因此
\begin{equation}
L_{\rm T}=\frac{4\pi\gamma_0^3m_{\rm e}v_0^3}{\hbar\omega^2}
=\SI{0.47770}{\metre}.
\label{eq:feqo-talbot-numerical}
\end{equation}
半复现长度 $L_{\rm T}/2$ 是把密度平移半周期的长度；原位复现需要上式的完整长度。
对 $b=1.20$，第一谐波的最早理想极大在 $L_1^*=\SI{0.02993}{\metre}$。

若源的中心能量有标准差 $\sigma_E=\SI{1}{\electronvolt}$ 的独立高斯涨落，
$\partial_Ev=1/(\gamma_0^3m_{\rm e}v_0)$ 给
\[
\sigma_t(L)=\frac{L\sigma_E}{\gamma_0^3m_{\rm e}v_0^3},\qquad
D_q(L)=\exp[-q^2\omega^2\sigma_t(L)^2/2].
\]
定义 $L_E$ 为第一谐波衰减到 $\ee^{-1/2}$ 的长度，则
\begin{equation}
L_E=\frac{\gamma_0^3m_{\rm e}v_0^3}{\omega\sigma_E}
=\SI{0.05891}{\metre}.
\label{eq:feqo-LE-def}
\end{equation}
在上述理想极大处 $D_1=0.879$，而 $D_2=0.597$；高次谐波先被洗掉。
因此最佳实际距离应最大化
$|J_1(4b\sin(2\pi L/L_{\rm T}))|\exp[-L^2/(2L_E^2)]$，
它通常比无噪声最优距离更短。

若改取 $K_0=\SI{100}{\kilo\electronvolt}$、$\lambda=\SI{2}{\micro\metre}$、$b=1$，
同一计算给 $L_{\rm T}=\SI{0.92869}{\metre}$、$L_1^*=\SI{0.07070}{\metre}$。
光学周期为 $\SI{6.6713}{\femto\second}$，但一个谐波极大本身不能定义脉冲全宽；
全宽必须从完整周期密度求峰值和半高交点，不能从 $b$ 单独指定一个通用常数。
明确一种可复查的定义：先选定一个峰 $\tau_*$ 和背景值 $n_{\rm bg}$，
在该峰相邻的两个单调区间解
\begin{equation}
 n(\tau_\pm)=n_{\rm bg}+\frac{n(\tau_*)-n_{\rm bg}}2,\qquad
 W_{1/2}=\tau_+-\tau_-.
 \label{eq:feqo-pulse-width-definition}
\end{equation}
绝对半高宽取 $n_{\rm bg}=0$；相对背景的半对比宽取实际选定背景。
弱聚束若整个周期密度都高于绝对半峰值，则绝对半高交点不存在。
多峰、周期包络或不同背景约定也会改变结果；完整密度以及这个读出定义应一起报告。
\end{exbox}


\subsection{Wigner--Moyal 输运}

在位置表象中容易看出波包有多宽，在动量表象中容易看出它含有哪些速度分量。若要理解一个带初始啁啾的波包为何先压缩再展宽，就需要同时保留位置与动量的相关。Wigner 表示把密度算符改写为位置 $z$ 和动量 $p$ 的函数，适合描述这种相关。它保留全部密度矩阵信息，但不必处处非负，因此需要从密度矩阵定义，而不能直接假设它是经典联合概率。

在纵向坐标表象中定义密度核
\[
\rho(z_1,z_2)\equiv\langle z_1|\hat\rho|z_2\rangle,
\qquad
\operatorname{Tr}\hat\rho=\int\dd z\,\rho(z,z)=1.
\]
对角元给位置概率，非对角元携带空间相干。为了同时保留平均位置与两点相干距离，引入中心坐标 $z=(z_1+z_2)/2$ 和相对坐标 $s=z_1-z_2$，并对 $s$ 做傅里叶变换：
\begin{equation}
\boxed{
W(z,p)
\equiv\frac{1}{2\pi\hbar}
\int_{-\infty}^{\infty}\dd s\,
\ee^{-\ii ps/\hbar}
\rho\!\left(z+\frac{s}{2},z-\frac{s}{2}\right).}
\label{eq:electron_wigner}
\end{equation}
其中 $W$ 的 SI 量纲是作用量$^{-1}$，使 $W\,\dd z\,\dd p$ 无量纲。它不是经典联合概率：相干交叉项可以使 $W$ 局部为负。但它的两个边缘严格恢复可测概率，
\begin{equation}
\boxed{
\int\dd p\,W(z,p)=\langle z|\hat\rho|z\rangle,
\qquad
\int\dd z\,W(z,p)=\langle p|\hat\rho|p\rangle.}
\label{eq:wigner_marginals}
\end{equation}
因此 Wigner 负值并不意味着“负概率事件”；它记录的是不能由正经典相空间分布表示的非对角相干。

波函数改用 Wigner 表示以后，Hamiltonian 和其他算符也需要相应的函数表示。对算符 $\hat A$，定义它的\emph{Weyl 符号}
\[
A_W(z,p)=\int\dd s\,\ee^{-\ii ps/\hbar}
\left\langle z+\frac{s}{2}\right|\hat A\left|z-\frac{s}{2}\right\rangle.
\]
注意此处不含 $1/(2\pi\hbar)$；密度算符的 Weyl 符号因此为 $\rho_W=2\pi\hbar W$。相应地，期待值为 $\langle\hat A\rangle=\int\dd z\,\dd p\,A_W W$。位置和动量算符的符号分别就是 $z$ 和 $p$。

两个算符的乘积一般不能用两个符号的普通乘积代替，因为 $\hat z\hat p$ 与 $\hat p\hat z$ 不同，而普通函数 $zp$ 与 $pz$ 相同。正确的对应乘法称为\emph{Moyal 星积}，定义为
\begin{equation}
\boxed{
(\hat A\hat B)_W=A_W\star B_W
=A_W\exp\!\left[
\frac{\ii\hbar}{2}
\left(
\overleftarrow\partial_z\overrightarrow\partial_p
-\overleftarrow\partial_p\overrightarrow\partial_z
\right)\right]B_W.}
\label{eq:weyl_star_product_derived}
\end{equation}
左箭头导数作用于左侧函数，右箭头导数作用于右侧函数。例如，$z\star p=zp+\ii\hbar/2$，而 $p\star z=pz-\ii\hbar/2$，相减正好恢复 $[\hat z,\hat p]=\ii\hbar$。指数中的高阶导数进一步保留一般算符乘法的量子修正。于是 Liouville--von Neumann 方程
\[
\ii\hbar\partial_t\hat\rho=[\hat H,\hat\rho]
\]
在相空间中精确等价于
\begin{equation}
\boxed{
\partial_tW=\{H_W,W\}_{\rm M}
\equiv\frac{1}{\ii\hbar}(H_W\star W-W\star H_W).}
\label{eq:moyal_general}
\end{equation}
当由 $\hbar$ 设置的相空间单元相对于 $H_W$ 与 $W$ 的变化尺度足够小时，星积展开的最低阶给出 Poisson 括号；经典 Liouville 输运因此是\eqnrefs{eq:moyal_general} 的受控长尺度极限。

\begin{theorem}[Weyl 乘积与 Moyal 演化]
对 Schwartz 符号 $A_W,B_W$，Weyl 量子化的乘积由\eqnrefs{eq:weyl_star_product_derived}的星积给出，指数记号按 Fourier 振荡积分定义；一般光滑符号的导数级数只作渐近解释。Liouville 方程相应变为\eqnrefs{eq:moyal_general}。二次多项式 Hamiltonian 的 Moyal 括号严格等于 Poisson 括号。
\end{theorem}
\begin{proof}
先从 Weyl 符号的 Fourier 表示出发。对足够良性的相空间函数 \(A_W(z,p)\)，写成
\[
A_W(z,p)
=
\int\frac{\dd\xi_z\,\dd\xi_p}{(2\pi)^2}\,
\widetilde A(\xi_z,\xi_p)
\ee^{\ii(\xi_z z+\xi_p p)},
\]
相应 Weyl 量子化为
\[
\hat A
=
\int\frac{\dd\xi_z\,\dd\xi_p}{(2\pi)^2}\,
\widetilde A(\xi_z,\xi_p)
\ee^{\ii(\xi_z\hat z+\xi_p\hat p)}.
\]
同理
\[
\hat B
=
\int\frac{\dd\xi_z'\,\dd\xi_p'}{(2\pi)^2}\,
\widetilde B(\xi_z',\xi_p')
\ee^{\ii(\xi_z'\hat z+\xi_p'\hat p)}.
\]

令
\[
X=\ii(\xi_z\hat z+\xi_p\hat p),
\qquad
Y=\ii(\xi_z'\hat z+\xi_p'\hat p).
\]
利用
\[
[\hat z,\hat p]=\ii\hbar
\]
得到
\[
[X,Y]
=
-\ii\hbar
(\xi_z\xi_p'-\xi_p\xi_z')\,I.
\]
该对易子是中心元，所以
\[
[X,[X,Y]]
=
[Y,[X,Y]]
=0,
\]
Baker--Campbell--Hausdorff 级数精确截断为
\[
\ee^X\ee^Y
=
\ee^{X+Y+\frac12[X,Y]}.
\]
于是
\begin{align*}
&\ee^{\ii(\xi_z\hat z+\xi_p\hat p)}
\ee^{\ii(\xi_z'\hat z+\xi_p'\hat p)}\\
&\qquad=
\exp\!\left[
-\frac{\ii\hbar}{2}
(\xi_z\xi_p'-\xi_p\xi_z')
\right]
\ee^{\ii[(\xi_z+\xi_z')\hat z+(\xi_p+\xi_p')\hat p]}.
\end{align*}

把这一乘法代回 \(\hat A\hat B\)，再取 Weyl 符号，得到
\begin{align*}
(\hat A\hat B)_W(z,p)
&=
\int
\frac{\dd\xi_z\,\dd\xi_p\,\dd\xi_z'\,\dd\xi_p'}{(2\pi)^4}\,
\widetilde A(\xi_z,\xi_p)
\widetilde B(\xi_z',\xi_p')\\
&\quad\times
\ee^{-\frac{\ii\hbar}{2}(\xi_z\xi_p'-\xi_p\xi_z')}
\ee^{\ii[(\xi_z+\xi_z')z+(\xi_p+\xi_p')p]}.
\end{align*}
另一方面，对 Fourier 基函数有
\[
\partial_z
\ee^{\ii(\xi_z z+\xi_p p)}
=
\ii\xi_z
\ee^{\ii(\xi_z z+\xi_p p)},
\qquad
\partial_p
\ee^{\ii(\xi_z z+\xi_p p)}
=
\ii\xi_p
\ee^{\ii(\xi_z z+\xi_p p)}.
\]
因此相位因子可以等价写成双向导数算符：
\[
\ee^{-\frac{\ii\hbar}{2}(\xi_z\xi_p'-\xi_p\xi_z')}
\longleftrightarrow
\exp\!\left[
\frac{\ii\hbar}{2}
\left(
\overleftarrow\partial_z\overrightarrow\partial_p
-
\overleftarrow\partial_p\overrightarrow\partial_z
\right)
\right].
\]
于是
\[
\boxed{
(\hat A\hat B)_W
=
A_W
\exp\!\left[
\frac{\ii\hbar}{2}
\left(
\overleftarrow\partial_z\overrightarrow\partial_p
-
\overleftarrow\partial_p\overrightarrow\partial_z
\right)
\right]
B_W
\equiv A_W\star B_W,
}
\]
即正文\eqnrefs{eq:weyl_star_product_derived}。

现在对 Liouville--von Neumann 方程
\[
\ii\hbar\partial_t\hat\rho
=
[\hat H,\hat\rho]
\]
取 Weyl 符号：
\[
\ii\hbar\partial_tW
=
H_W\star W-W\star H_W.
\]
两边除以 \(\ii\hbar\)，得到
\[
\boxed{
\partial_tW
=
\{H_W,W\}_{\rm M}
\equiv
\frac1{\ii\hbar}
(H_W\star W-W\star H_W),
}
\]
即正文\eqnrefs{eq:moyal_general}。

若把星积按 \(\hbar\) 展开，偶数阶在交换子中相消，只保留奇数阶：
\begin{align*}
\{H_W,W\}_{\rm M}
&=
\left(
\partial_zH_W\,\partial_pW
-\partial_pH_W\,\partial_zW
\right)\\
&\quad
-\frac{\hbar^2}{24}
\left(
\partial_z^3H_W\,\partial_p^3W
-3\partial_z^2\partial_pH_W\,
  \partial_p^2\partial_zW
+3\partial_z\partial_p^2H_W\,
  \partial_p\partial_z^2W
-\partial_p^3H_W\,\partial_z^3W
\right)
+\mathcal O(\hbar^4).
\end{align*}
最低阶正是 Poisson 括号，
\[
\{H_W,W\}_{\rm M}
=
\{H_W,W\}_{\rm PB}
+\mathcal O(\hbar^2).
\]
\end{proof}


对精确自由相对论电子，哈密顿量只依赖动量，
\[
H_W(p)=E(p)=\sqrt{m_{\mathrm e}^2c^4+c^2p^2},
\qquad
v(p)=E'(p)=\frac{c^2p}{E(p)}.
\]
Moyal 方程于是化成奇数空间导数级数
\begin{equation}
\boxed{
\partial_tW
=-E'(p)\partial_zW
+\frac{\hbar^2}{24}E'''(p)\partial_z^3W
-\frac{\hbar^4}{1920}E^{(5)}(p)\partial_z^5W+\cdots.}
\label{eq:relativistic_free_moyal}
\end{equation}
第一项就是沿精确群速度 $v(p)$ 的经典平流；其余项衡量“非二次色散 + 有限量子相空间分辨率”造成的量子修正。若只保留中心动量 $p_0$ 附近的二阶色散
\[
E(p_0+\delta p)
\simeq E_0+v_0\delta p+\frac{\delta p^2}{2m_\parallel},
\qquad
m_\parallel\equiv\frac{1}{E''(p_0)}=\gamma_0^3m_{\mathrm e},
\]
所有三阶以上动量导数都消失，所以 Moyal 方程在这个二次模型内\emph{精确}退化成相空间剪切：
\begin{equation}
\boxed{
\partial_tW+
\left(v_0+\frac{p-p_0}{m_\parallel}\right)\partial_zW=0,
\qquad
W(z,p,t)=
W_0\!\left(z-\left[v_0+\frac{p-p_0}{m_\parallel}\right]t,p\right).}
\label{eq:wigner_shear}
\end{equation}

自由传播不改变每个动量分量本身，却让不同动量以不同群速度前进。于是相空间分布沿 $z$ 方向剪切；所谓“电子啁啾压缩”就是先制备一张倾斜的相空间分布，再让这张分布被自由剪切到位置投影最窄的时刻。

\begin{proposition}[二次色散的精确 Wigner 剪切]
在 $E_2(p)=E_0+v_0(p-p_0)+(p-p_0)^2/(2m_\parallel)$ 模型中，Wigner 演化严格为\eqnrefs{eq:wigner_shear}；完整色散的高阶导数不能据此置零。对 Schwartz 初态此式逐点成立，对一般态按分布解释。
\end{proposition}
\begin{proof}
自由电子 Hamiltonian 只依赖动量，$H_W=E(p)$。从 Moyal 方程的星积定义出发，位置导数只能作用在 $W$ 上，因此
\begin{align*}
 \partial_tW
 &=\frac{1}{\ii\hbar}(E\star W-W\star E)\\
 &=-\sum_{n=0}^{\infty}
 \frac{(-1)^n}{(2n+1)!}
 \left(\frac{\hbar}{2}\right)^{2n}
 E^{(2n+1)}(p)\,\partial_z^{2n+1}W,
\end{align*}
这就是正文\eqnrefs{eq:relativistic_free_moyal}。现在在中心动量 $p_0$ 附近定义二次模型
\begin{equation*}
 E_2(p)=E_0+v_0(p-p_0)+\frac{(p-p_0)^2}{2m_\parallel}.
\end{equation*}
因为 $E_2^{(r)}(p)=0$ 对所有 $r\ge3$ 成立，Moyal 级数不是“近似丢掉量子项”，而是在这个二次 Hamiltonian 内严格终止：
\begin{equation*}
 \partial_tW+
 \left(v_0+\frac{p-p_0}{m_\parallel}\right)\partial_zW=0.
\end{equation*}
这是一个一阶线性偏微分方程。令参数 $s$ 沿特征线变化，要求
\begin{equation*}
 \frac{\dd}{\dd s}W[z(s),p(s),t(s)]
 =\partial_tW\,\dot t+
 \partial_zW\,\dot z+
 \partial_pW\,\dot p=0.
\end{equation*}
与偏微分方程逐项比较，可选
\begin{equation*}
 \dot t=1,
 \qquad
 \dot p=0,
 \qquad
 \dot z=v_0+\frac{p-p_0}{m_\parallel}.
\end{equation*}
于是 $p(s)=p(0)$ 保持不变，而
\begin{equation*}
 z(t)=z(0)+
 \left(v_0+\frac{p-p_0}{m_\parallel}\right)t.
\end{equation*}
沿同一条特征线 $W$ 为常数，因此把初始位置反解为
$z(0)=z-[v_0+(p-p_0)/m_\parallel]t$，得到
\begin{equation*}
 W(z,p,t)=W_0\!\left(
 z-\left[v_0+\frac{p-p_0}{m_\parallel}\right]t,
 p\right),
\end{equation*}
即正文\eqnrefs{eq:wigner_shear}。

“剪切”一词现在有精确几何含义：同一 $p$ 水平线上的分布只做刚性平移，但平移距离依赖 $p$，所以原本竖直的相空间线元会倾斜。若使用完整相对论 $E(p)$，三阶及更高导数重新出现，演化不再是这个线性辛剪切；其误差必须由正文给出的高阶累计相位条件控制。
\end{proof}


把这个剪切写成矩阵，就可以直接计算位置、动量方差以及它们的协方差怎样变化。

在随中心电子移动的坐标
\[
\zeta\equiv z-v_0t,
\qquad
\delta p\equiv p-p_0
\]
中，二阶自由哈密顿量的线性漂移已经被消去，留下
\[
H_{\rm com}=\frac{\delta p^2}{2m_\parallel}.
\]
因此 Heisenberg 演化为
\begin{equation}
\boxed{
\zeta(t)=\zeta(0)+\frac{t}{m_\parallel}\delta p(0),
\qquad
\delta p(t)=\delta p(0).}
\label{eq:electron-shear-heisenberg}
\end{equation}
这与\eqnrefs{eq:wigner_shear} 是同一剪切在算符与相空间中的两种写法。

定义中心化二阶矩
\[
\sigma_z^2\equiv\langle\zeta^2\rangle,
\qquad
\sigma_p^2\equiv\langle\delta p^2\rangle,
\qquad
C_{zp}\equiv\frac12\langle\zeta\delta p+\delta p\zeta\rangle.
\]
\eqnrefs{eq:electron-shear-heisenberg} 给出
\begin{equation}
\boxed{
\sigma_z^2(t)
=\sigma_z^2(0)+\frac{2t}{m_\parallel}C_{zp}(0)
+\frac{t^2}{m_\parallel^2}\sigma_p^2.}
\label{eq:electron_shear_variance}
\end{equation}
负协方差 $C_{zp}<0$ 表示“前后位置与动量反相关”，自由剪切会先减小位置宽度。最佳时间与最小宽度为
\begin{equation}
\boxed{
 t_*=-m_\parallel\frac{C_{zp}(0)}{\sigma_p^2},
\qquad
\sigma_{z,\min}^2
=\sigma_z^2(0)-\frac{C_{zp}^2(0)}{\sigma_p^2}
\ge\frac{\hbar^2}{4\sigma_p^2}.}
\label{eq:electron_chirp_compression_bound}
\end{equation}
所以啁啾不会压缩相空间面积；它只是把已有的相空间椭圆剪切到位置投影最窄。$m_\parallel$ 决定焦点出现的传播时间，却不改变固定 $\sigma_p$ 下的 Heisenberg 极限。

\begin{proposition}[二次剪切的方差极小值]
设正则共轭对具有有限二阶矩，协方差使用中心化的对称乘积，$\sigma_p^2>0$。二次自由演化满足\eqnrefs{eq:electron_shear_variance}；在全实时间上，$t_*=-m_\parallel C_{zp}/\sigma_p^2$，且 $\sigma_{z,\min}^2=\sigma_z^2-C_{zp}^2/\sigma_p^2\ge\hbar^2/(4\sigma_p^2)$。若只许 $t\ge0$，则需 $t_*\ge0$，否则最小值在 $t=0$。
\end{proposition}
\begin{proof}
设初态二阶矩有限，不要求波包为 Gaussian。共动二次 Hamiltonian 为
\[
H_{\rm com}=\frac{\delta p^2}{2m_\parallel},
\qquad
[\zeta,\delta p]=\ii\hbar.
\]
先在 Heisenberg 绘景计算算符运动。因为
\begin{align*}
[\zeta,\delta p^2]
&=[\zeta,\delta p]\,\delta p
+\delta p\,[\zeta,\delta p]\\
&=2\ii\hbar\,\delta p,
\end{align*}
所以
\begin{align*}
\dot\zeta
&=\frac{1}{\ii\hbar}[\zeta,H_{\rm com}]
=\frac{\delta p}{m_\parallel},\\
\dot{\delta p}
&=\frac{1}{\ii\hbar}[\delta p,H_{\rm com}]=0.
\end{align*}
积分得到
\[
\zeta(t)=\zeta(0)+\frac{t}{m_\parallel}\delta p(0),
\qquad
\delta p(t)=\delta p(0).
\]
由于 $[\zeta(0),\delta p(0)]=\ii\hbar$，平方时必须保留交叉项的算符次序：
\begin{align*}
\zeta(t)^2
&=\zeta^2(0)
+\frac{t}{m_\parallel}
\bigl[\zeta(0)\delta p(0)+\delta p(0)\zeta(0)\bigr]
+\frac{t^2}{m_\parallel^2}\delta p^2(0).
\end{align*}
将上面的算符恒等式用于中心化变量 $\Delta\zeta=\zeta-\langle\zeta\rangle$、$\Delta p=\delta p-\langle\delta p\rangle$ 后取期望值，并使用对称协方差定义
\[
C_{zp}(0)=\frac12
\langle\Delta\zeta(0)\Delta p(0)+\Delta p(0)\Delta\zeta(0)\rangle,
\]
便得到
\[
\sigma_z^2(t)
=\sigma_z^2(0)
+\frac{2t}{m_\parallel}C_{zp}(0)
+\frac{t^2}{m_\parallel^2}\sigma_p^2,
\]
即正文\eqnrefs{eq:electron_shear_variance}。把右端视为关于 $t$ 的二次函数，求导得
\[
\frac{\dd\sigma_z^2(t)}{\dd t}
=\frac{2}{m_\parallel}C_{zp}(0)
+\frac{2t}{m_\parallel^2}\sigma_p^2.
\]
若 $\sigma_p^2>0$，驻点唯一且二阶导数
\[
\frac{\dd^2\sigma_z^2(t)}{\dd t^2}
=\frac{2\sigma_p^2}{m_\parallel^2}>0,
\]
因此它确为最小值。解得
\[
t_*=-m_\parallel\frac{C_{zp}(0)}{\sigma_p^2}.
\]
代回二次函数：
\begin{align*}
\sigma_{z,\min}^2
&=\sigma_z^2(0)
-2\frac{C_{zp}^2(0)}{\sigma_p^2}
+\frac{C_{zp}^2(0)}{\sigma_p^2}\\
&=\sigma_z^2(0)-\frac{C_{zp}^2(0)}{\sigma_p^2}.
\end{align*}
最后对正则共轭对 $(\zeta,\delta p)$ 使用 Robertson--Schr\"odinger 不等式
\[
\sigma_z^2(0)\sigma_p^2-C_{zp}^2(0)
\ge\frac14|\langle[\zeta,\delta p]\rangle|^2
=\frac{\hbar^2}{4}.
\]
由于 $\sigma_p^2>0$，两边除以 $\sigma_p^2$ 得
\[
\sigma_z^2(0)-\frac{C_{zp}^2(0)}{\sigma_p^2}
\ge\frac{\hbar^2}{4\sigma_p^2}.
\]
与上式合并即得到正文\eqnrefs{eq:electron_chirp_compression_bound}。
\end{proof}


无初始 $z$--$p$ 相关的最小不确定度高斯波包是最简单的基准。若初始纵向宽度为 $\sigma_{z0}$，则二阶传播保持高斯形状，并有
\[
\sigma_z^2(t)
=\sigma_{z0}^2+
\left(\frac{\hbar t}{2m_\parallel\sigma_{z0}}\right)^2.
\]
横向二次哈密顿量为 $H_\perp=(\delta p_x^2+\delta p_y^2)/(2m_\perp)$。若取其中任一横向方向 $x$，初态同样为无相关最小不确定度高斯，宽度 $\sigma_{x0}$ 与动量宽度满足 $\sigma_{p_x}=\hbar/(2\sigma_{x0})$，则
\[
\sigma_x^2(t)
=\sigma_{x0}^2+
\left(\frac{\hbar t}{2m_\perp\sigma_{x0}}\right)^2.
\]
定义展宽时间为使第二项与初始方差相等的时间，
\[
t_{\rm disp,\parallel}
=\frac{2m_\parallel\sigma_{z0}^2}{\hbar},
\qquad
t_{\rm disp,\perp}
=\frac{2m_\perp\sigma_{x0}^2}{\hbar}.
\]
对相同初始空间宽度 $\sigma_{x0}=\sigma_{z0}$，
\[
\frac{t_{\rm disp,\parallel}}{t_{\rm disp,\perp}}
=\frac{m_\parallel}{m_\perp}=\gamma_0^2.
\]
对高斯初态，前述二阶色散可以解析求解，并直接给出展宽与压缩条件。

二阶模型的适用域由精确色散中未被 $E_2$ 吸收的余项决定。令 $\delta p=p-p_0$。Taylor 定理给出精确分解
\[
E(p_0+\delta p)
=E_0+v_0\delta p+\frac{\delta p^2}{2m_\parallel}
+\frac{E'''_0}{6}\delta p^3+R_4(\delta p),
\qquad
E'''_0=-\frac{3v_0}{\gamma_0^4m_{\mathrm e}^2c^2},
\]
其中积分余项定义为
\[
R_4(\delta p)
=\frac{\delta p^4}{6}
\int_0^1\dd s\,(1-s)^3
E^{(4)}(p_0+s\delta p).
\]
对一标准差工作窗 $|\delta p|\le\sigma_p$ 定义
\[
\epsilon_3(t)
\equiv\frac{|E'''_0|t\sigma_p^3}{6\hbar},
\qquad
\epsilon_4(t)
\equiv
\frac{t\sigma_p^4}{24\hbar}
\sup_{|p-p_0|\le\sigma_p}|E^{(4)}(p)|.
\]
于是精确色散截断量满足
\begin{equation}
\boxed{
\epsilon_{\rm disp}^{(\ge3)}(t;\sigma_p)
\equiv
\sup_{|\delta p|\le\sigma_p}
\left|\frac{t}{\hbar}[E(p_0+\delta p)-E_2(p_0+\delta p)]\right|
\le\epsilon_3(t)+\epsilon_4(t).}
\label{eq:electron_dispersion_cubic_control}
\end{equation}
因此 $\epsilon_3$ 只有在同时满足 $\epsilon_4\ll\epsilon_3$ 时才可单独作为“三阶主导”的简化指标；二阶传播本身的受控条件是右端总上界远小于 $1$。一旦这一条件失效，应回到精确传播相位 $\exp[-\ii E(p)t/\hbar]$，而不是引入经验“修正质量”。

\begin{proposition}[相对论色散的三阶相位控制]
设动量严格支撑于 $|p-p_0|\le\sigma_p$，$t\ge0$。二次 Taylor 传播的相位余项不超过 $|E^{(3)}_0|t\sigma_p^3/(6\hbar)+t\sigma_p^4\sup|E^{(4)}|/(24\hbar)$，即\eqnrefs{eq:electron_dispersion_cubic_control}。仅有均方带宽约束时，应另外保留窗外概率误差。
\end{proposition}
\begin{proof}
取局部工作动量窗 $\mathcal W_{p,1}=\{p:|p-p_0|\le\sigma_p\}$，令 $\delta p=p-p_0$。自由电子的精确正能色散为
\[
E(p)=\sqrt{m_{\mathrm e}^2c^4+c^2p^2}.
\]
为了把余项界写成可直接计算的量，先从这个一般表达式求高阶导数。一次求导给
\[
E'(p)=\frac{c^2p}{E(p)}.
\]
再次求导时保留商的两项：
\begin{align*}
E''(p)
&=\frac{c^2}{E(p)}
-\frac{c^2p}{E^2(p)}E'(p)\\
&=\frac{c^2}{E(p)}-\frac{c^4p^2}{E^3(p)}\\
&=\frac{m_{\mathrm e}^2c^6}{E^3(p)}.
\end{align*}
因此
\begin{align*}
E'''(p)
&=-3m_{\mathrm e}^2c^6E^{-4}(p)E'(p)\\
&=-\frac{3m_{\mathrm e}^2c^8p}{E^5(p)},
\end{align*}
在 $p=p_0$ 处恢复正文使用的
\[
E'''_0=-\frac{3v_0}{\gamma_0^4m_{\mathrm e}^2c^2}.
\]
继续求一次导数可得四阶项
\begin{align*}
E^{(4)}(p)
&=-3m_{\mathrm e}^2c^8
\left[E^{-5}(p)-5pE^{-6}(p)E'(p)\right]\\
&=-\frac{3m_{\mathrm e}^2c^8}{E^5(p)}
+\frac{15m_{\mathrm e}^2c^{10}p^2}{E^7(p)}\\
&=\frac{3m_{\mathrm e}^2c^8}{E^7(p)}
\left(4c^2p^2-m_{\mathrm e}^2c^4\right).
\end{align*}
这些导数决定二次传播模型被舍弃的谱相位尺度。

对 $E(p_0+\delta p)$ 在 $p_0$ 处使用三阶 Taylor 公式的积分余项，
\begin{align*}
E(p_0+\delta p)-E_2(p_0+\delta p)
&=\frac{E'''_0}{6}\delta p^3+R_4(\delta p),\\
R_4(\delta p)
&=\frac{\delta p^4}{6}
\int_0^1\dd s\,(1-s)^3
E^{(4)}(p_0+s\delta p).
\end{align*}
积分余项保留了从 $p_0$ 到 $p_0+\delta p$ 整条线段上的四阶导数。因为 $|\delta p|\le\sigma_p$ 且 $0\le s\le1$，积分中的动量始终满足 $|p_0+s\delta p-p_0|\le\sigma_p$。于是
\begin{align*}
|R_4(\delta p)|
&\le\frac{|\delta p|^4}{6}
\int_0^1\dd s\,(1-s)^3
\left|E^{(4)}(p_0+s\delta p)\right|\\
&\le\frac{|\delta p|^4}{6}
\sup_{|p-p_0|\le\sigma_p}|E^{(4)}(p)|
\int_0^1\dd s\,(1-s)^3\\
&=\frac{|\delta p|^4}{24}
\sup_{|p-p_0|\le\sigma_p}|E^{(4)}(p)|\\
&\le\frac{\sigma_p^4}{24}
\sup_{|p-p_0|\le\sigma_p}|E^{(4)}(p)|.
\end{align*}
把这个结果代入正文\eqnrefs{eq:dispersion_truncation_control} 的定义，得到
\begin{align*}
\epsilon_{\rm disp}^{(\ge3)}(t;\sigma_p)
&=\frac{t}{\hbar}
\sup_{|\delta p|\le\sigma_p}
\left|E(p_0+\delta p)-E_2(p_0+\delta p)\right|\\
&\le
\frac{t}{\hbar}
\sup_{|\delta p|\le\sigma_p}
\left(\frac{|E'''_0|}{6}|\delta p|^3+|R_4(\delta p)|\right)\\
&\le
\frac{|E'''_0|t\sigma_p^3}{6\hbar}
+\frac{t\sigma_p^4}{24\hbar}
\sup_{|p-p_0|\le\sigma_p}|E^{(4)}(p)|\\
&=\epsilon_3(t)+\epsilon_4(t),
\end{align*}
即正文\eqnrefs{eq:electron_dispersion_cubic_control}。
\end{proof}


对 PINEM/QPINEM 产生的离散相干能量梳，同一剪切可在能量--时间傅里叶对偶中读取。设边带能量为 $E_0+\ell\hbar\omega$，并把随中心电子运动的延迟写成
\[
\tau\equiv t-\frac{z}{v_0}.
\]
若共同慢包络为 $F(\tau)$，边带振幅为 $c_\ell$，则
\begin{equation}
\boxed{
\Psi(\tau)=F(\tau)\sum_\ell c_\ell\ee^{-\ii\ell\omega\tau},
\qquad
I(\tau)=|F(\tau)|^2
\sum_{\ell,m}c_\ell c_m^*\ee^{-\ii(\ell-m)\omega\tau}.}
\label{eq:coherent_energy_comb_time_amplitude}
\end{equation}
只有非对角相干 $c_\ell c_m^*$ 才产生光学周期或亚周期密度结构；边带概率很宽却彼此无相位关系时，不会自动形成阿秒脉冲列。若相干支持大致延伸到 $|\ell|\le L$，其能量带宽 $\mathcal B_E^{\rm coh}\sim2L\hbar\omega$ 给出傅里叶尺度
\[
\Delta t\gtrsim\frac{\hbar}{\mathcal B_E^{\rm coh}},
\]
但真正达到这一尺度还要求传播后的二次及更高阶谱相位得到控制。

在固定传播距离 $L_z$ 上，最方便的是展开纵向动量 $p(E)$。由 $\dd p/\dd E=1/v$ 和 $\dd v/\dd p=1/m_\parallel$，中心能量处
\[
\left.\frac{\dd^2 p}{\dd E^2}\right|_{E_0}
=-\frac{1}{m_\parallel v_0^3}.
\]
因此第 $\ell$ 个边带在去掉共同相位与线性时间平移后获得二次谱相位
\begin{equation}
\boxed{
\varphi_\ell^{(2)}(L_z)
=-\ell^2\vartheta_2(L_z),
\qquad
\vartheta_2(L_z)
\equiv\frac{\hbar\omega^2L_z}{2m_\parallel v_0^3}.}
\label{eq:energy_comb_quadratic_phase}
\end{equation}
由 $\dd p/\dd E=1/v$ 与 $\dd v/\dd p=1/m_\parallel$ 逐步展开，二阶 Taylor 系数因此完全由中心群速度和纵向色散质量确定。
因此 $\vartheta_2(L_z)$ 直接衡量离散能量格的二次传播相位，并决定时间聚焦与 Talbot 相位匹配条件。

\begin{figure}[htbp]
\centering
\includegraphics[width=0.76\linewidth]{gaussian_compression.png}
\caption{带负 $z$--$p$ 协方差的高斯电子包在二阶相对论色散下先压缩后展宽。焦点由\eqnrefs{eq:electron_chirp_compression_bound} 决定。}
\label{fig:gaussian-compression}
\end{figure}

Wigner 边缘性质属于定义层面的 Fourier 代数，不另设推导框。对\eqnrefs{eq:electron_wigner} 积分动量并用

\[
\int\dd p\,\ee^{-\ii ps/\hbar}=2\pi\hbar\,\delta(s),
\]
得到
\[
\int\dd p\,W(z,p)
=\rho(z,z).
\]
反过来积分 $z$，作变量变换 $z_1=z+s/2,z_2=z-s/2$，雅可比行列式为一，并使用全文约定
\[
\langle z|p\rangle=(2\pi\hbar)^{-1/2}\ee^{\ii pz/\hbar},
\]
于是
\begin{align*}
\int\dd z\,W(z,p)
&=\frac{1}{2\pi\hbar}
\int\dd z_1\dd z_2\,
\ee^{-\ii p(z_1-z_2)/\hbar}\rho(z_1,z_2)\\
&=\langle p|\hat\rho|p\rangle.
\end{align*}
这同时固定了\eqnrefs{eq:electron_wigner} 前因子 $1/(2\pi\hbar)$。

$\epsilon_3$ 衡量二次模型遗漏的首个谱相位，$\epsilon_4$ 控制继续截断在三阶本身的误差；一般应检查两者之和，而不能只看三阶项。

经典 Liouville 方程是 Moyal 演化的最低阶梯度极限；对至多二次的 Hamiltonian，三阶及更高导数严格为零，因此量子 Wigner 演化与经典相空间平流精确重合。

\subsection{时间透镜与 Talbot 传播}

前面的相空间计算表明，若先让波包不同位置具有不同平均速度，自由传播就可能使它压缩。对纵向电子波包，随到达时间变化的光学相位可以产生这种速度相关，作用类似空间光学中的薄透镜，因此称为\emph{时间透镜}。另一方面，若波包含有等间隔的相干能量边带，各边带的二次传播相位还可能在特定距离重新对齐，使周期时间图样复现；这称为\emph{Talbot 复现}。下面分别由传播相位计算聚焦条件和复现条件。

对经典相干光场，PINEM 章节定义的无量纲耦合 $\beta=|\beta|\ee^{\ii\varphi_\beta}$ 决定电子获得的周期相位。用到达延迟 $\tau=t-z/v_0=-\zeta/v_0$ 表示，能量阶数 $\ell$ 的 Fourier 因子固定为 $\ee^{-\ii\ell\omega\tau}$，因此，单次相位调制写成
\begin{equation}
\boxed{
\Psi(\tau)\longrightarrow
\exp\!\left[2\ii|\beta|\sin(\varphi_\beta-\omega\tau)\right]\Psi(\tau).}
\label{eq:electron_periodic_time_phase}
\end{equation}
它不是全局抛物线，因此“时间透镜”只在选定光学相位附近成立。若以相位极值 $\varphi_\beta-\omega\tau_0=\pi/2$ 为中心，令 $\delta\tau=\tau-\tau_0$，则局域展开给
\begin{equation}
\boxed{
\varphi_{\rm lens}(\delta\tau)
=-|\beta|\omega^2\delta\tau^2
+O\!\left(|\beta|\omega^4\delta\tau^4\right).}
\label{eq:electron_local_time_lens_phase}
\end{equation}
常数相位已省略。控制“抛物线透镜”近似的参数是 $|\omega\delta\tau|\ll1$；超出这一时间窗，正弦相位的高阶项会生成周期性焦散/脉冲列，而不是单个理想高斯焦点。

\begin{proposition}[周期相位的局部二次展开]
设周期相位在指定中心附近三次可微。在半宽 $a$ 的时间窗内，用其二次 Taylor 多项式作局部透镜相位时，点态乘法传播子的误差至多为相位余项的上确界，后者不超过 $a^3\sup|\phi^{(3)}|/6$。本章正弦调制的中心相位、群延迟和二次曲率由\eqnrefs{eq:electron_local_time_lens_phase}给出。
\end{proposition}
\begin{proof}
在相位极值附近对正弦调制做 Taylor 展开，推出\eqnrefs{eq:electron_local_time_lens_phase} 并识别四阶项作为有限孔径误差。 透镜近似要求 $|\omega\delta\tau|\ll1$； 超出这一局域时间窗时应使用完整的正弦相位调制。 定义光学相位

\begin{equation*}
\theta(\tau)=\varphi_\beta-\omega\tau .
\end{equation*}
电子的周期相位调制为 $2|\beta|\sin\theta$。选择相位极值
\begin{equation*}
\theta_0=\frac{\pi}{2},
\qquad
\varphi_\beta-\omega\tau_0=\frac{\pi}{2},
\end{equation*}
并写
\begin{equation*}
\delta\tau=\tau-\tau_0,
\qquad
\delta\theta=-\omega\delta\tau.
\end{equation*}
利用
\begin{equation*}
\sin\!\left(\frac\pi2+\delta\theta\right)
=\cos\delta\theta
=1-\frac{\delta\theta^2}{2}
+\frac{\delta\theta^4}{24}
+O(\delta\theta^6),
\end{equation*}
相位指数的实函数变成
\begin{align*}
2|\beta|\sin\theta
&=2|\beta|
-|\beta|\delta\theta^2
+\frac{|\beta|}{12}\delta\theta^4
+O(|\beta|\delta\theta^6)\\
&=2|\beta|
-|\beta|\omega^2\delta\tau^2
+\frac{|\beta|\omega^4\delta\tau^4}{12}
+O(|\beta|\omega^6\delta\tau^6).
\end{align*}
第一项只产生与 $\delta\tau$ 无关的整体相位，去掉它以后得到正文的局域二次相位
\begin{equation*}
\varphi_{\rm lens}(\delta\tau)
=-|\beta|\omega^2\delta\tau^2
+O(|\beta|\omega^4\delta\tau^4).
\end{equation*}
四次项相对于二次项的绝对比值为
\begin{equation*}
\left|
\frac{(|\beta|/12)\omega^4\delta\tau^4}
{|\beta|\omega^2\delta\tau^2}
\right|
=\frac{(\omega\delta\tau)^2}{12}.
\end{equation*}
\end{proof}


局域二次时间相位先建立到达时间与能量的相关；随后自由色散把这种相关转化为实际的到达时间压缩。粗粒化地看，\eqnrefs{eq:classical_energy_mod_from_beta} 给相位 $\theta=\varphi_\beta-\omega\tau$ 上的能量调制
\[
\delta E(\theta)=2\hbar\omega|\beta|\cos\theta.
\]
经过距离 $L_z$ 后，速度色散把它映射成到达时间偏移。固定传播距离时
\begin{align*}
\tau(E)&=\frac{L_z}{v(E)},\\
\left.\frac{\dd(1/v)}{\dd E}\right|_{E_0}
&=-\frac1{v_0^2}
\left.\frac{\dd v/\dd p}{\dd E/\dd p}\right|_{p_0}
=-\frac{1}{m_\parallel v_0^3},
\end{align*}
其中用了 $\dd E/\dd p=v$ 与 $\dd v/\dd p=1/m_\parallel$。因此对小能量调制
\begin{equation}
\boxed{
\delta\tau_{\rm drift}(\theta)
=
-\frac{2L_z\hbar\omega|\beta|}{m_\parallel v_0^3}\cos\theta
+O\!\left([\delta E(\theta)]^2\right).}
\label{eq:electron_classical_bunching_map}
\end{equation}
这里的一阶式只要求在实际占据能宽内 $1/v(E)$ 的二阶及更高变化可忽略；若能量调制大到越过这一局域线性区，应直接使用精确飞行时间 $L_z/v(E_0+\delta E)$。当映射 $\tau_{\rm out}(\theta)$ 的雅可比行列式变小时，经典束流图像中出现密度聚束；量子图像中，同一现象由\eqnrefs{eq:coherent_energy_comb_time_amplitude} 的边带交叉项干涉形成。两种语言在边带不可分辨、包络跨越许多光学周期的粗粒化极限相接。

另一种特殊相位匹配不是局域聚焦，而是整个离散能量梳的谱相位复现。\secref{subsec:sideband-to-temporal-density-dp46}已经从完整密度谐波得到唯一的 Talbot 时间和长度~\eqnrefs{eq:electron_talbot_time_length_dp46}。这里用能量表象检验同一结果：由\eqnrefs{eq:energy_comb_quadratic_phase}，当 $\vartheta_2(L_T)=2\pi$ 时，任意整数 $\ell$ 都满足 $\exp[-\ii\ell^2\vartheta_2(L_T)]=1$；代入 $\vartheta_2=\hbar\omega^2L_z/(2m_\parallel v_0^3)$ 恰好恢复\eqnrefs{eq:electron_talbot_time_length_dp46} 在 $\kappa_z=\omega/v_0$ 下的 $L_T$。在 $L_z/L_T$ 为有理数时，$\ell^2$ 的模整数结构产生分数 Talbot 图样。初始能散、有限包络、空间电荷以及\eqnrefs{eq:electron_dispersion_cubic_control} 所衡量的高阶色散都会降低复现对比度。

\begin{figure}[H]
  \centering
  \includegraphics[width=0.82\linewidth]{talbot_time_focusing.png}
  \caption{周期相位调制后的自由电子在二次谱相位传播下形成的时域 Talbot 结构。图中从 $c_\ell=J_\ell(2|\beta|)$ 出发，对每个边带乘以 $\exp(-\ii s\ell^2)$ 后重新求和；取 $|\beta|=3$。\cite{Feist2015Nature}}
  \label{fig:talbot_time_focusing}
\end{figure}

经典纵向束流动力学使用 $(z,\delta E)$ 相空间描述啁啾与微聚束。\cite{Pellegrini2016FEL} 当单电子相干、反冲和负 Wigner 结构不可忽略时，Wigner 表示保留同一几何直觉并补上量子相干。因此二者不是互斥理论，而是不同分辨尺度上的连续描述。

精确相对论色散决定电子闭合传播，相空间表示把同一动力学改写成几何剪切，周期相位与二次谱相位则决定聚焦和复现。未追踪环境会把闭合幺正动力学变成约化量子信道，而探测器进一步把量子态映射为经典数据；二者分别由开放系统动力学与量子测量描述。

上述复现只涉及二次色散。若占据带宽内的三阶及更高谱相位在一个 $L_T$ 上积累到 $O(1)$，不同整数边带不能同时回相；这一偏离由正文\eqnrefs{eq:electron-talbot-higher-order-control-dp76} 的 $\epsilon_T^{(\ge3)}$ 控制。

抛物线时间透镜的局域控制量是 $|\omega\delta\tau|$；它不再远小于一时，应使用完整周期正弦相位。

\subsection{有限高斯波包：解析解 TDSE 与时间聚焦}
\label{sec:electron-gaussian-lens-analytic}

时间透镜的聚焦也可以直接从一个有限波包的波函数计算。取已建立的纵向二阶色散，在随中心电子运动的坐标 $\zeta=z-v_0t$ 中，包络满足含时 Schrödinger 方程（time-dependent Schrödinger equation, TDSE）
\begin{equation}
 \ii\hbar\partial_t\psi(\zeta,t)
 =-\frac{\hbar^2}{2m_\parallel}\partial_\zeta^2\psi(\zeta,t),
 \qquad \int|\psi|^2\dd\zeta=1.
 \label{eq:electron-gaussian-envelope-tdse}
\end{equation}
这个方程与自由非相对论粒子的 Schrödinger 方程形式相同，但系数 $m_\parallel=\gamma_0^3m_{\mathrm e}$ 来自相对论能量的二次展开。中心能量和匀速运动已被移入载波相位，$\psi$ 只描述相对于中心电子的包络变化。它的适用条件仍是实际动量范围内高阶色散积累的相位可以忽略。

用复参数 $a_0$ 表示一般纯高斯输入，$\Re a_0>0$：
\begin{equation*}
 \psi(\zeta,0)=\mathcal N_0\ee^{-a_0\zeta^2},\qquad
 \mathcal N_0=\left(\frac{2\Re a_0}{\pi}\right)^{1/4}.
\end{equation*}
$a_0$ 的实部决定密度宽度，虚部决定二次空间相位；$[a_0]=\mathrm{m^{-2}}$。连续 TDSE 的精确高斯解为
\begin{equation}
 \psi(\zeta,t)=\frac{\mathcal N_0}{\sqrt{d(t)}}
 \exp\!\left[-\frac{a_0\zeta^2}{d(t)}\right],\qquad
 d(t)=1+\frac{2\ii\hbar a_0t}{m_\parallel}.
 \label{eq:electron-gaussian-tdse-solution}
\end{equation}
平方根支路从 $d(0)=1$ 连续选取。$d$ 无量纲；前因子同时包含保持归一化所需的幅度变化与共同传播相位，不能只保留指数中的宽度变化。

\begin{exbox}[带二次相位的波包何时最窄]
令初始密度标准差为 $\sigma_{\zeta0}$，空间啁啾参数为实数 $C_\zeta$，即
\begin{equation*}
 a_0=\frac{1-\ii C_\zeta}{4\sigma_{\zeta0}^2},\qquad
 t_{\rm dif}=\frac{2m_\parallel\sigma_{\zeta0}^2}{\hbar},\qquad
 u_t=\frac{t}{t_{\rm dif}}.
\end{equation*}
这里 $C_\zeta$ 是无量纲实数，控制位置表象中的二次相位；前面用 $C_p$ 表示的是动量表象中的二次相位。$t_{\rm dif}$ 具有时间单位，它是无啁啾高斯包络因自由传播而明显展宽的时间尺度，$u_t$ 则以这个时间为单位。代入一般解，有 $d=1+(C_\zeta+\ii)u_t$。利用
\begin{equation*}
 \Re\frac{a_0}{d}=\frac{\Re a_0}{|d|^2},\qquad
 |d|^2=(1+C_\zeta u_t)^2+u_t^2,
\end{equation*}
便从概率密度的高斯指数读出
\begin{equation*}
 \sigma_\zeta^2(t)=\sigma_{\zeta0}^2
 \left[1+2C_\zeta u_t+(1+C_\zeta^2)u_t^2\right].
\end{equation*}
这是关于时间的二次式，配方即可求极小值：
\begin{equation*}
 \frac{\sigma_\zeta^2(t)}{\sigma_{\zeta0}^2}
 =(1+C_\zeta^2)
 \left(u_t+\frac{C_\zeta}{1+C_\zeta^2}\right)^2
 +\frac1{1+C_\zeta^2}.
\end{equation*}
若 $C_\zeta<0$，未来存在焦点，
\begin{equation}
 t_f=-\frac{C_\zeta}{1+C_\zeta^2}t_{\rm dif},\qquad
 L_f=v_0t_f,\qquad
 \sigma_{\zeta,\min}=\frac{\sigma_{\zeta0}}{\sqrt{1+C_\zeta^2}}.
 \label{eq:electron-gaussian-focus}
\end{equation}
若 $C_\zeta=0$，波包从初始时刻起展宽；若 $C_\zeta>0$，二次式的焦点位于过去，未来同样单调展宽。因此聚焦需要正确符号的相位相关，而不只是足够大的能宽。

若输入本来无啁啾，在时间透镜的相位极值附近乘入 $\exp[-\ii|\beta|\omega^2\tau^2]$，并使用 $\zeta=-v_0\tau$，便得到
\begin{equation*}
 C_\zeta=-4|\beta|\omega^2\sigma_{\tau0}^2,
 \qquad \sigma_{\tau0}=\sigma_{\zeta0}/v_0.
\end{equation*}
其符号恰好允许未来聚焦；将它代入\eqnrefs{eq:electron-gaussian-focus} 就得到对任意允许参数成立的焦距与压缩比。到达时间宽度取 $\sigma_\tau\simeq\sigma_\zeta/v_0$，还要求速度展宽相对中心速度很小。
\end{exbox}

为了理解压缩过程中动量分布怎样变化，再计算位置与动量的协方差。记 $\Delta\zeta=\hat\zeta-\langle\hat\zeta\rangle$、$\Delta p=\hat p-\langle\hat p\rangle$；由于两个算符不对易，用对称乘积定义实协方差。对上述初态，在位置表象使用 $\hat p=-\ii\hbar\partial_\zeta$，得
\begin{equation*}
 \sigma_p^2=\frac{\hbar^2(1+C_\zeta^2)}{4\sigma_{\zeta0}^2},\qquad
 C_{\zeta p}\equiv\tfrac12\langle\Delta\zeta\Delta p+\Delta p\Delta\zeta\rangle
 =\frac{\hbar C_\zeta}{2}.
\end{equation*}
所以 $\sigma_{\zeta0}^2\sigma_p^2-C_{\zeta p}^2=\hbar^2/4$。自由传播保持 $\sigma_p$ 不变，只把协方差变为 $C_{\zeta p}(t)=C_{\zeta p}(0)+t\sigma_p^2/m_\parallel$；焦点恰在该协方差过零时，且 $\sigma_{\zeta,\min}\sigma_p=\hbar/2$。因此焦点处位置与动量不再相关，宽度乘积恰好达到最小不确定值。

时间透镜的适用条件也应作用于\emph{相位本身}。若检查包络的 $|\tau|\le K\sigma_{\tau0}$ 区域，其中 $K$ 是选定的覆盖倍数，则正弦相位余项的首项为
\begin{equation*}
 |\delta\varphi_4|\lesssim\frac{|\beta|}{12}
 (\omega K\sigma_{\tau0})^4.
\end{equation*}
要可靠地用二次相位替代完整作用，应同时检查 $\omega K\sigma_{\tau0}\ll1$ 和 $|\delta\varphi_4|\ll1$，并检查高阶色散在 $t_f$ 内的累积误差。仅有局部几何孔径很小，在极强驱动下仍不足以保证波函数误差很小。



\begin{derivation}[高斯波包的自由传播与方程核对]
在二阶包络模型中，连续动量谱给传播核
\begin{equation*}
 K(\zeta,\zeta';t)=\int\frac{\dd p}{2\pi\hbar}
 \exp\!\left[\frac{\ii p(\zeta-\zeta')}{\hbar}
             -\frac{\ii p^2t}{2m_\parallel\hbar}\right]
 =\sqrt{\frac{m_\parallel}{2\pi\ii\hbar t}}
 \exp\!\left[\frac{\ii m_\parallel(\zeta-\zeta')^2}{2\hbar t}\right],
 \quad t>0.
\end{equation*}
振荡高斯积分可先加正的收敛因子，再取其趋零极限。将初态代入 $\psi(\zeta,t)=\int K\psi(\zeta',0)\dd\zeta'$，整理为
\begin{equation*}
 -A\zeta'^2+B\zeta'+C,
 \quad A=a_0-\frac{\ii m_\parallel}{2\hbar t},\quad
 B=-\frac{\ii m_\parallel\zeta}{\hbar t},\quad
 C=\frac{\ii m_\parallel\zeta^2}{2\hbar t}.
\end{equation*}
$\Re A>0$ 保证积分收敛；使用 $\int\ee^{-Ax^2+Bx}\dd x=\sqrt{\pi/A}\ee^{B^2/(4A)}$，合并前因子与指数即得\eqnrefs{eq:electron-gaussian-tdse-solution}。

还可独立代回检验。写 $a(t)=a_0/d(t)$、$\mathcal N(t)=\mathcal N_0/\sqrt{d(t)}$，则
\begin{equation*}
 \dot a=-\frac{2\ii\hbar}{m_\parallel}a^2,\qquad
 \frac{\dot{\mathcal N}}{\mathcal N}=-\frac{\ii\hbar}{m_\parallel}a,
 \qquad
 \frac{\partial_\zeta^2\psi}{\psi}=-2a+4a^2\zeta^2.
\end{equation*}
含时 Schrödinger 方程左右两侧的常数项与 $\zeta^2$ 项于是分别相等。再用
$\Re(a_0/d)=\Re a_0/|d|^2$，对 $|\psi|^2$ 作高斯积分，归一化保持为一。所以传播后的函数同时满足初始条件、运动方程和归一化条件。
\end{derivation}

